% ==============================================================================
% ch01.tex -- Chapter 1: Generalities on "A-Structures" and "A-Generators"
% ==============================================================================
\chapter{Generalities on ``$A$-Structures'' and ``$A$-Generators''}
\label{ch:1}

% ------------------------------------------------------------------------------
% Section 1.1
% ------------------------------------------------------------------------------
\section{Review of relative Cartier divisors (Compare \cite{Mum2}, pp.~61--73.)}
\label{sec:1.1}

\kmitem\label{item:1.1.1}
Let $S$ be an arbitrary scheme, and let $X$ be an $S$-scheme. By an \emph{effective Cartier divisor $D$ in $X/S$} we mean a closed subscheme $D \subset X$ such that
\begin{equation*}
\begin{cases}
D \text{ is flat over } S, \\
\text{the ideal sheaf } \cI(D) \subset \cO_X \text{ is an invertible } \cO_X\text{-module, i.e.,} \\
\text{it is a locally free } \cO_X\text{-module of rank one.}
\end{cases}
\end{equation*}
This notion is local on $S$. When $S$ is affine, say $S = \Spec(R)$, it means that we can cover $X$ by affine opens $U_i = \Spec(A_i)$, $A_i$ an $R$-algebra, such that $D \cap U_i$ is defined in $U_i$ by one equation $f_i = 0$, where $f_i \in A_i$ is an element such that
\begin{equation*}
\begin{cases}
A_i/f_i A_i \text{ is flat over } R, \\
f_i \text{ is not a zero-divisor in } A_i.
\end{cases}
\end{equation*}
The tautological exact sequence on $X$,
\[
0 \longrightarrow \cI(D) \longrightarrow \cO_X \longrightarrow \cO_D \longrightarrow 0,
\]
becomes on $U_i = \Spec(A_i)$ the exact sequence
\[
0 \longrightarrow A_i \xrightarrow{\;\times f_i\;} A_i \longrightarrow A_i/f_i A_i \longrightarrow 0.
\]

\kmitem\label{item:1.1.2}
Given two effective Cartier divisors $D$ and $D'$ in $X/S$, their \emph{sum} $D + D'$ is the effective Cartier divisor in $X/S$ defined locally by the product of the defining equations of $D$ and $D'$. Explicitly, if $S = \Spec(R)$ and if on an affine open $\Spec(A)$ of $X$, $D$ and $D'$ are defined respectively by equations $f = 0$ and $g = 0$, then $D + D'$ is defined there by $fg = 0$. To check that $fg$ is not a zero-divisor in $A$, one notes the commutative diagram
\[
\begin{tikzcd}[column sep=huge]
A \arrow[r, "\times f"] \arrow[rr, bend right=30, "\times fg"'] & A \arrow[r, "\times g"] & A.
\end{tikzcd}
\]
To check that $A/fgA$ is flat over $R$, one notes the short exact sequence
\[
0 \longrightarrow A/gA \xrightarrow{\;\times f\;} A/fgA \longrightarrow A/fA \longrightarrow 0,
\]
which exhibits $A/fgA$ as an extension of flat $R$-modules.

\kmitem\label{item:1.1.3}
Given an effective Cartier divisor $D$ in $X/S$, we may speak of the inverse (as invertible $\cO_X$-module) $\cI^{-1}(D)$ of its ideal sheaf. We have a tautological exact sequence
\[
0 \longrightarrow \cO_X \longrightarrow \cI^{-1}(D) \longrightarrow \cO_D \otimes_{\cO_X} \cI^{-1}(D) \longrightarrow 0.
\]
The inclusion of $\cO_X$ in $\cI^{-1}(D)$ allows us to view the constant function ``$1$'' as a global section of $\cI^{-1}(D)$, and we may recover $D$ as the scheme of zeroes of this global section of $\cI^{-1}(D)$.

Conversely, suppose we are given a pair $(\cL, \ell)$ consisting of an invertible $\cO_X$-module $\cL$ on $X$ together with a global section $\ell \in \mathrm{H}^0(X, \cL)$ which sits in a short exact sequence of $\cO_X$-modules
\[
0 \longrightarrow \cO_X \xrightarrow{\;\times \ell\;} \cL \longrightarrow \cL/\cO_X \longrightarrow 0
\]
with $\cL/\cO_X$ flat over $S$. Then the scheme of zeroes of the section $\ell$ of $\cL$ is easily seen to be an effective Cartier divisor $D$ in $X/S$, and there is a unique isomorphism of $(\cL, \ell)$ with $(\cI^{-1}(D), \text{``}1\text{''})$.

This construction allows us to interpret effective Cartier divisors in $X/S$ as isomorphism classes of pairs $(\cL, \ell)$ as above. From this point of view, the operation ``sum of effective Cartier divisors in $X/S$'' is none other than the operation of tensor product:
\[
(\cL, \ell) + (\cL', \ell') = (\cL \otimes \cL', \ell \otimes \ell').
\]
The zero element for this addition is the pair $(\cO_X, 1)$, corresponding to the empty Cartier divisor.

\kmitem\label{item:1.1.4}
There are two natural situations in which one can define the inverse image of a relative Cartier divisor. First let
\[
T \longrightarrow S
\]
be an arbitrary morphism of schemes. Then for any effective Cartier divisor $D$ in $X/S$, say represented by a pair $(\cL, \ell)$, the closed subscheme $D_T \stackrel{\mathrm{dfn}}{=} D \times_S T$ of $X_T = X \times_S T$ is an effective Cartier divisor in $X_T/T$, represented by the pair $(\cL_T, \ell_T)$ on $X_T$. To see this, it suffices to treat the case when $S = \Spec(R)$ and $T = \Spec(R')$ are both affine; then the sequence on $X_T = X \otimes_R R'$
\[
\cO_{X \otimes_R R'} \xrightarrow{\;\ell \otimes 1\;} \cL \otimes_R R' \longrightarrow (\cL \otimes_R R')/(\cO_{X \otimes_R R'})
\]
is obtained from the short exact sequence on $X$
\[
0 \longrightarrow \cO_X \xrightarrow{\;\ell\;} \cL \longrightarrow \cL/\cO_X \longrightarrow 0
\]
by applying the functor $\otimes_R R'$. Because $\cL/\cO_X$ is assumed flat over $R$, this sequence stays short exact after $\otimes_R R'$, and its last term is $R'$-flat. Therefore $(\cL \otimes_R R', \ell \otimes 1)$ defines an effective Cartier divisor in $X \otimes_R R'/R'$, as required.

Second, let
\[
\begin{tikzcd}
Y \arrow[r, "f"] \arrow[rd] & X \arrow[d] \\
& S
\end{tikzcd}
\]
be a flat morphism of $S$-schemes. Then any effective Cartier divisor $D$ in $X/S$ gives rise to an effective Cartier divisor $f^*(D)$ in $Y/S$. Indeed, the cartesian diagram
\[
\begin{tikzcd}
f^*(D) \arrow[r, hook] \arrow[d] & Y \arrow[d, "f\text{ flat}"] \\
D \arrow[r, hook] & X
\end{tikzcd}
\]
shows that $f^*(D)$ is flat over $D$, and hence, $D$ being $S$-flat, that $f^*(D)$ is flat over $S$. To see that the ideal sheaf $\cI(f^*(D))$ is an invertible $\cO_Y$-module, we remark that this ideal sheaf is none other than $f^*(\cI(D))$, as follows from the fact that the short exact sequence on $X$
\[
0 \longrightarrow \cI(D) \longrightarrow \cO_X \longrightarrow \cO_D \longrightarrow 0
\]
remains short exact after application of the functor $f^*$, thanks to the flatness of $f$.

\kmitem\label{item:1.1.5}
We now turn to the question of recognizing which closed subschemes of $X$ are in fact effective Cartier divisors in $X/S$.

\begin{customprop}{1.1.5.1}\label{prop:1.1.5.1}
Suppose that $S$ is locally noetherian, and that $X$ is an $S$-scheme of finite type which is flat over $S$. Let $\cF$ be a coherent sheaf on $X$ which is flat over $S$. Then the necessary and sufficient condition that $\cF$ be flat over $\cO_X$ is that for every geometric point of $S$, i.e., every morphism $\Spec(k) \to S$ with $k$ an algebraically closed field, the induced coherent sheaf $\cF \otimes_R k$ on the fibre $X \otimes_R k$ be flat over $\cO_{X \otimes_R k}$.
\end{customprop}

\begin{proof}
This is just the fibre-by-fibre criterion of flatness \cite[V, 3.6]{AK1}.
\end{proof}

\begin{customcor}{1.1.5.2}\label{cor:1.1.5.2}
Let $S$ be locally noetherian, $X$ a flat $S$-scheme of finite type, and $D \subset X$ a closed subscheme which is flat over $S$. Then $D$ is an effective Cartier divisor in $X/S$ if and only if for all geometric points $\Spec(k) \to S$ of $S$, the closed subscheme $D \otimes_R k$ of $X \otimes_R k$ is an effective Cartier divisor in $X \otimes_R k / k$.
\end{customcor}

\begin{proof}
The necessity is a special case of the preservation of effective Cartier divisors in $X/S$ under arbitrary change of base $T \to S$. For the sufficiency, we apply the proposition to $\cF = \cI(D)$, which is $S$-flat because it sits in the short exact sequence
\[
0 \longrightarrow \cI(D) \longrightarrow \cO_X \longrightarrow \cO_D \longrightarrow 0,
\]
in which both $\cO_X$ and $\cO_D$ are $S$-flat by hypothesis. Tensoring over $\cO_S$ with $k$ (for any geometric point $\Spec(k) \to S$), this sequence stays exact ($S$-flatness of $\cO_D$), so comparing first terms yields
\[
\cI(D) \otimes_R k \xrightarrow{\;\sim\;} \cI(D \otimes_R k).
\]
By hypothesis, $\cI(D \otimes k)$ is an invertible $\cO_{X \otimes k}$-module, so in particular it is $\cO_{X \otimes k}$-flat. Therefore $\cI(D)$ is $\cO_X$-flat by the proposition. Because $\cI(D)$ is also coherent, it is a locally free $\cO_X$-module. That it is locally free of rank one may then be seen by restricting it to the fibres $X \otimes k$ of $X/S$, on each of which it is locally free of rank one.
\end{proof}

% ------------------------------------------------------------------------------
% Section 1.2
% ------------------------------------------------------------------------------
\section{Relative Cartier divisors in curves}
\label{sec:1.2}

\kmitem\label{item:1.2.1}
Let $S$ be an arbitrary scheme. By a \emph{smooth curve $C/S$}, we will always mean a smooth morphism $C \to S$ of relative dimension one which is separated and of finite presentation.

\begin{lemma}\label{lem:1.2.2}
If $C/S$ is a smooth curve, then any section $s \in C(S)$ defines an effective Cartier divisor, noted $[s]$, in $C/S$.
\end{lemma}

\begin{proof}
Because $C/S$ is of finite presentation, we are immediately reduced to the case when $S = \Spec(R)$ with $R$ noetherian (\cite[IV, 8.9.1]{EGA}), then by Corollary~\ref{cor:1.1.5.2} to the case when $R$ is an algebraically closed field $k$, and in this case the assertion is obvious.
\end{proof}

\begin{lemma}\label{lem:1.2.3}
Let $C/S$ be a smooth curve as above. Let $D \subset C$ be a closed subscheme which is finite and flat over $S$, and of finite presentation over $S$. Then $D$ is an effective Cartier divisor in $C/S$, which is proper over $S$. Conversely every effective Cartier divisor in $C/S$ which is proper over $S$ is of this form.
\end{lemma}

\begin{proof}
To prove the first statement, we are immediately reduced (by \cite[IV, 8.9.1 and 11.2.6.1]{EGA}) to the case when $S = \Spec(R)$ with $R$ noetherian, then by Corollary~\ref{cor:1.1.5.2} to the case when $R$ is an algebraically closed field $k$, in which case it is obvious that $D$ is an effective Cartier divisor. That $D$ is proper over $S$ follows from its being assumed finite over $S$.

Conversely, any effective Cartier divisor $D$ in $C/S$ is certainly of finite presentation, so we are reduced to the case when $S = \Spec(R)$ with $R$ noetherian (\cite[IV, 8.9.1 and 11.2.6.1]{EGA}). We must prove that $D$ is finite over $R$. But $D$ is proper over $R$, so it suffices to prove that $D/R$ has finite fibres. Thus we are reduced to the case when $R$ is an algebraically closed field $k$, i.e., to the assertion that an effective Cartier divisor in a smooth curve over $k$ is a finite $k$-scheme, which is obvious.
\end{proof}

\begin{remark}\label{rem:1.2.4}
If $C/S$ is a proper smooth curve, then every effective Cartier divisor in $C/S$ is automatically proper over $S$.
\end{remark}

\begin{definition}\label{def:1.2.5}
Let $C/S$ be a smooth curve, $D$ an effective Cartier divisor in $C/S$ which is proper over $S$. Then locally on $S$, say $S = \Spec(R)$, the affine ring of $D$ is a locally free $R$-module of finite rank. This rank, which is constant Zariski locally on $S$, is called the degree of $D$, noted $\deg(D)$. In terms of an $(\cL, \ell)$ presentation of $D$, the exact sequence on $C$
\[
0 \longrightarrow \cO_C \xrightarrow{\ \ell\ } \cL \longrightarrow \cL/\cO \longrightarrow 0
\]
has
\[
\resizebox{\linewidth}{!}{$\displaystyle
\cL/\cO \cong \cI^{-1}(D)/\cO \cong \cI^{-1}(D) \tens{\cO_C} \cO_D \cong \cL \tens{\cO_C} \cO_D = \text{an invertible } \cO_D\text{-module}.
$}
\]
Therefore $D$ is proper over $S$ if and only if the sheaf $\cL/\cO$ on $C$ has its support proper over $S$. If this is the case, then locally on $S$, say $S = \Spec(R)$, the $R$-module
\[
\mathrm{H}^0(C, \cL/\cO) = \mathrm{H}^0(C, \cI^{-1}(D)/\cO)
\]
is a locally free $R$-module of rank $=$ degree$(D)$.
\end{definition}

\begin{lemma}\label{lem:1.2.6}
Let $C/S$ be a smooth curve, and let $D_1, D_2$ be two effective Cartier divisors in $C/S$ which are both proper over $S$. Then $D_1 + D_2$ is proper over $S$, and we have the equality
\[
\deg(D_1 + D_2) = \deg(D_1) + \deg(D_2).
\]
\end{lemma}

\begin{proof}
Let $(\cL_1, \ell_1)$ and $(\cL_2, \ell_2)$ represent $D_1$ and $D_2$. The short exact sequence on $C$
\[
0 \longrightarrow \cL_2/\cO \longrightarrow \cL_1 \otimes \cL_2/\cO \longrightarrow \cL_1 \otimes \cL_2/\cL_2 \longrightarrow 0
\]
is obtained by applying the snake lemma to the diagram
\[
\begin{tikzcd}[column sep=small, every label/.append style={font=\scriptsize}]
0 \arrow[r] & \cO \arrow[r, "\times 1"] \arrow[d, "\times \ell_2"'] & \cO \arrow[r] \arrow[d, "\times \ell_1 \otimes \ell_2"] & 0 \arrow[r] \arrow[d, "0"] & 0 \\
0 \arrow[r] & \cL_2 \arrow[r, "\times \ell_1"'] & \cL_1 \otimes \cL_2 \arrow[r] & \cL_1 \otimes \cL_2 / \cL_2 \arrow[r] & 0
\end{tikzcd}
\]
when the bottom row is obtained by applying $\otimes_{\cO} \cL_2$ to
\[
0 \longrightarrow \cO \xrightarrow{\;\times \ell_1\;} \cL_1 \longrightarrow \cL_1/\cO \longrightarrow 0.
\]
In any case, the above exact sequence, when rewritten
\[
\begin{tikzcd}[row sep=tiny, column sep=tiny, every label/.append style={font=\scriptsize}]
0 \arrow[r] & \cL_2 \tens{\cO} \cO_{D_2} \arrow[r] & (\cL_1 \otimes \cL_2) \tens{\cO} \cO_{D_1+D_2} \arrow[r] & (\cL_1 \otimes \cL_2) \tens{\cO} \cO_{D_1} \arrow[r] & 0 \\
& \rotatebox{90}{$=$} & \rotatebox{90}{$=$} & \rotatebox{90}{$=$} & \\
& (\cL_2)|_{D_2} & (\cL_1 \otimes \cL_2)|_{D_1+D_2} & (\cL_1 \otimes \cL_2)|_{D_1} &
\end{tikzcd}
\]
shows that $D_1 + D_2$ is proper over $S$ and shows, taking Euler characteristics, that the degrees add as asserted.
\end{proof}

\begin{lemma}\label{lem:1.2.7}
Let $C/S$ be a smooth curve. Then for any section $s \in C(S)$, the associated effective Cartier divisor $[s]$ in $C/S$ is proper over $S$, and has degree one. Conversely, any effective Cartier divisor $D$ in $C/S$ which is proper over $S$ and has degree one is of the form $[s]$ for a unique section $s \in C(S)$.
\end{lemma}

\begin{proof}
If $s \in C(S)$, then the Cartier divisor $[s]$ sits in a diagram
\[
\begin{tikzcd}
{[s]} \arrow[r, hook] \arrow[rd, "\sim"'] & C \arrow[d] \\
& S,
\end{tikzcd}
\]
which shows that $[s]$ is proper over $S$, of degree one.

Conversely, if $D$ is an effective Cartier divisor in $C/S$, proper over $S$ of degree one, then we have a diagram
\[
\begin{tikzcd}
D \arrow[r, hook] \arrow[rd] & C \arrow[d] \\
& S
\end{tikzcd}
\]
in which the diagonal arrow is an isomorphism (because locally on $S$, say $S = \Spec(R)$, it turns the affine ring of $D$ into an $R$-algebra which is an invertible $R$-module, i.e., into $R$ itself).
\end{proof}

\begin{lemma}\label{lem:1.2.8}
Let $C', C$ be two smooth curves over $S$, and
\[
\begin{tikzcd}
C' \arrow[r, "f"] \arrow[rd] & C \arrow[d] \\
& S
\end{tikzcd}
\]
a morphism which is finite and flat of degree noted $\deg(f)$. If $D$ is an effective Cartier divisor in $C/S$, proper over $S$, then its inverse image $f^*(D)$ in $C'/S$ is an effective Cartier divisor in $C'/S$, proper over $S$, and its degree is given by
\[
\deg(f^*(D)) = \deg(f) \deg(D).
\]
\end{lemma}

\begin{proof}
That $f^*(D)$ is an effective Cartier divisor in $C'/S$ we have already seen in \ref{item:1.1.4}. That it is proper over $S$ and that its degree is given by the stated formula is immediate from the fact that $\cO_{f^*(D)} = f^*(\cO_D)$ as $\cO_{C'}$-module, so that $\cO_{f^*(D)}$ as $\cO_S$-module is given by $f_*(\cO_{f^*(D)}) = f_*(f^*(\cO_D)) = f_*(\cO_{C'}) \otimes_{\cO_C} \cO_D$.
\end{proof}

\begin{lemma}\label{lem:1.2.9}
Let $C/S$ be a smooth curve, $D \subset C$ an effective Cartier divisor in $C/S$ which is proper over $S$. Let $T \to S$ be an arbitrary morphism of schemes. Then the inverse image $D_T = D \times_S T$ is an effective Cartier divisor in $C_T/T$, proper over $T$, and of the same degree as $D$:
\[
\deg(D_T) = \deg(D).
\]
\end{lemma}

\begin{proof}
That $D_T$ is an effective Cartier divisor in $C_T/T$ we have seen in \ref{item:1.1.4}. That it is proper over $T$ and of the same degree as $D$ follows from the fact that $\cO_{D_T} = \cO_D \otimes_{\cO_S} \cO_T$.
\end{proof}

% ------------------------------------------------------------------------------
% Section 1.3
% ------------------------------------------------------------------------------
\section{Existence of incidence schemes}
\label{sec:1.3}

\kmitem\label{item:1.3.1}
Given two effective Cartier divisors $D, D'$ in $X/S$, we say that $D' \le D$ if $D' \subset D$ (inclusion of closed subschemes of $X$), i.e., if $\cI(D) \subset \cI(D')$ (inclusion of ideal sheaves inside $\cO_X$), or equivalently if there exists an effective Cartier divisor $D''$ in $X/S$ such that
\[
D = D' + D''.
\]
In terms of local equations $f = 0$ for $D$ and $g = 0$ for $D'$, the condition $D' \le D$ means precisely that $g$ divides $f$, i.e., that locally we have $f = gh$ for some function $h$. This function $h$ is unique (because $g$ is not a zero-divisor), and $h$ itself is not a zero divisor (because $gh = f$ is not a zero divisor). The locally defined subschemes $h = 0$ of $X$ patch together to define a subscheme $D''$ of $X$, whose ideal sheaf $\cI(D'')$, locally $h\cO_X$, is invertible. To see that $D''$ is flat over $S$, we may assume $S = \Spec(R)$, $X = \Spec(A)$, $D$ is defined by $f = 0$, $D'$ by $g = 0$, and $D''$ by $h = 0$, with $f = gh$. The exact sequence
\[
0 \longrightarrow A/hA \xrightarrow{\;\times g\;} A/fA \longrightarrow A/gA \longrightarrow 0
\]
then shows that $A/hA$ is $R$-flat (because $A/fA$ and $A/gA$ are $R$-flat), as required.

Globally, in terms of representatives $(\cL, \ell)$ for $D$ and $(\cL', \ell')$ for $D'$, the condition $D' \le D$ is precisely that the global section $\ell$ of $\cL$ vanish identically in $\cL|_{D'} = \cL \otimes_{\cO_X} \cO_{D'}$. Then $D''$ is represented by $(\cL'', \ell'')$ with $\cL'' = \cL \otimes (\cL')^{-1}$ and $\ell'' = \ell/\ell'$. The global version of the above short exact sequence is
\begin{equation}\label{eq:1.3.1.1}
0 \longrightarrow \cL'' \otimes_{\cO_X} \cO_{D''} \longrightarrow \cL \otimes_{\cO_X} \cO_D \longrightarrow \cL \otimes_{\cO_X} \cO_{D'} \longrightarrow 0.
\end{equation}

\kmitem\label{item:1.3.2}
We now turn from a general $X/S$ to a smooth curve $C/S$ (cf.~\ref{item:1.2.1}).

\begin{lemma}\label{lem:1.3.3}
Let $C/S$ be a smooth curve, $D$ and $D'$ effective Cartier divisors in $C/S$ which are both proper over $S$. If $D' \le D$, then the effective Cartier divisor $D'' = D - D'$ is proper over $S$, and
\[
\deg(D'') = \deg(D) - \deg(D').
\]
\end{lemma}

\begin{proof}
The properness is obvious from the short exact sequence \eqref{eq:1.3.1.1}, since both $D$ and $D'$ are proper over $S$. The degree assertion then results from \ref{lem:1.2.6}.
\end{proof}

\begin{keylemma}\label{lem:1.3.4}
Let $C/S$ be a smooth curve, $D$ and $D'$ effective Cartier divisors in $C/S$, with $D'$ proper over $S$. Then
\begin{enumerate}[(1)]
\item there exists a unique closed subscheme $Z \subset S$ which is universal for the relation $D' \le D$ in the following sense: given any morphism of schemes $T \to S$, the inverse images $D'_T$ and $D_T$ in $C_T$ satisfy $D'_T \le D_T$ if and only if the morphism $T \to S$ factors through $Z$;
\item the subscheme $Z \subset S$ is defined locally on $S$ by $\deg(D')$ equations;
\item formation of the closed subscheme $Z \subset S$ commutes with arbitrary change of base $S' \to S$, in the sense that the closed subscheme $Z'$ of $S'$ ``universal for the relation $D'_{S'} \le D_{S'}$'' is none other than $Z \times_S S'$.
\end{enumerate}
\end{keylemma}

\begin{proof}
The question is clearly local on $S$, which we may assume affine, say $S = \Spec(R)$. In terms of a representative $(\cL, \ell)$ for $D$, the condition $D' \le D$ is that the global section $\ell$ of $\cL$ vanish identically in $\cL \otimes_{\cO_C} \cO_{D'} = \cL|_{D'}$. Because $D'$ is finite locally-free over $S$, and $\cL|_{D'}$ is an invertible $\cO_{D'}$-module, the module $\mathrm{H}^0(D', \cL|_{D'})$ is a locally free $R$-module of $\text{rank} = \deg(D')$. Locally on $R$, we may choose an $R$-basis $e_1, \dots, e_{d'}$ of this $R$-module. The element $\ell$ has a unique expression
\[
\ell = \sum_{i=1}^{\deg(D')} r_i e_i, \quad \text{coefficients } r_i \in R.
\]
The condition ``$\ell = 0$'' is then represented by the closed subscheme of $\Spec(R)$ defined by the simultaneous vanishing of $r_1, \dots, r_{d'}$.
\end{proof}

\begin{corollary}\label{cor:1.3.5}
Let $C/S$ be a smooth curve, $D$ and $D'$ effective Cartier divisors in $C/S$ which are both proper over $S$ of the same degree. Then there exists a unique closed subscheme $Z \subset S$ which is universal for the relation $D = D'$, in the sense that for any $T \to S$, we have $D_T = D'_T$ in $C_T$ if and only if $T \to S$ factors through $Z$. The subscheme $Z \subset S$ is defined locally in $S$ by $\deg(D)$ equations, and its formation commutes with arbitrary change of base $S' \to S$.
\end{corollary}

\begin{proof}
The condition $D = D'$ is equivalent to the condition $D' \le D$, because $D$ and $D'$ both have the same degree.
\end{proof}

\kmitem\label{item:1.3.6}
Let $C/S$ be a smooth curve which is a group-scheme over $S$. Let $D$ be an effective Cartier divisor in $C/S$ which is proper over $S$. We say that $D$ is a \emph{subgroup} of $C/S$ if for every $S$-scheme $T$ the subset $D(T)$ of the group $C(T)$ is in fact a subgroup. This amounts to the existence of a (necessarily unique) structure of finite flat $S$-group-scheme on $D$ such that $D \hookrightarrow C$ is an $S$-group-scheme homomorphism.

\begin{corollary}\label{cor:1.3.7}
Let $C/S$ be a smooth curve which is a group-scheme over $S$. Let $D$ be an effective Cartier divisor in $C/S$, proper over $S$. Then there exists a unique closed subscheme $Z \subset S$ which is universal for the relation ``$D$ is a subgroup,'' in the sense that for any $T \to S$, $D_T$ in $C_T/T$ is a subgroup if and only if $T \to S$ factors through $Z$. The subscheme $Z \subset S$ is defined locally in $S$ by $1 + \deg(D) + (\deg(D))^2$ equations, and its formation commutes with arbitrary change of base $S' \to S$.
\end{corollary}

\begin{proof}
Let us denote by $e \in C(S)$ the identity section for the group structure, by
\[
\mathrm{inv} \colon C \longrightarrow C
\]
the $S$-automorphism of $C$ expressing inversion, and by
\[
m \colon C \underset{S}{\times} C \longrightarrow C
\]
the $S$-morphism ``multiplication.''

In order that $D$ be a subgroup, it is necessary and sufficient that the following three conditions be satisfied:
\begin{enumerate}[(1)]
\item Denoting by $[e]$ the effective Cartier divisor in $C/S$ attached to the identity section $e \in C(S)$, we have the inequality of Cartier divisors in $C/S$
\[
[e] \le D.
\]
(This is the condition that the subset $D(T) \subset C(T)$ always contain the identity.)
\item We have the equality of effective Cartier divisors in $C/S$
\[
D = \mathrm{inv}^*(D).
\]
(This is the condition that the subset $D(T) \subset C(T)$ be stable by inversion.)
\item Let $W$ denote the $S$-scheme $D \underset{S}{\times} D$, which represents the functor ``ordered pairs of points of $D$.'' Let $(P_1, P_2)$ be the universal pair of points of $D$, so that we have a tautological diagram
\[
\begin{tikzcd}[column sep=huge, row sep=large]
D_W \arrow[r, hook] & C_W \arrow[d] & D_W \arrow[l, hook'] \\
& W \arrow[lu, bend left=25, "P_1"] \arrow[ru, bend right=25, "P_2"'] &
\end{tikzcd}
\]
Then we must have the inequality of effective Cartier divisors in $C_W$
\[
[m(P_1, P_2)] \le D_W.
\]
(This is the condition that the subset $D(T) \subset C(T)$ always be stable by multiplication.)
\end{enumerate}

Now condition (1) is represented locally on $S$ by one equation. Condition (2) is represented locally on $S$ by $\deg(D)$ equations. Condition (3) is represented locally on $W$ by one equation, but because $W$ is finite locally free over $S$ of rank $(\deg(D))^2$, the vanishing on $W$ of a single function is equivalent to the vanishing on $S$ of its $(\deg(D))^2$ ``coordinates.''
\end{proof}

% ------------------------------------------------------------------------------
% Section 1.4
% ------------------------------------------------------------------------------
\section{Points of ``exact order $N$'' and cyclic subgroups}
\label{sec:1.4}

\kmitem\label{item:1.4.1}
Let $C/S$ be a smooth curve (cf.~\ref{item:1.2.1}), given with a structure of commutative $S$-group-scheme. We will refer to such a $C/S$ as a ``smooth commutative $S$-group-scheme of relative dimension one''; in view of \ref{item:1.2.1}, $C/S$ is automatically separated and of finite presentation over $S$. Let $N \ge 1$ be an integer. We say that a point $P \in C(S)$ has ``exact order $N$'' if the effective Cartier divisor $D$ in $C/S$ of degree $N$ defined by
\begin{equation}\label{eq:1.4.1.1}
D \xlongequal{dfn} [P] + [2P] + \cdots + [NP]
\end{equation}
is a \emph{subgroup} of $C/S$. We call this subgroup the cyclic subgroup of rank $N$ ``generated'' by $P$. We say that a closed subgroup-scheme $G \subset C$ which is finite locally-free of rank $N$ over $S$ is \emph{cyclic} if, locally f.p.p.f.\ \cite[III, Exp.~IV, 6.3]{SGA} on $S$, $G$ admits a generator, i.e., if there exists some f.p.l.p.f.\ (faithfully flat, locally of finite presentation) morphism $T \to S$ and a point $P \in C(T) = C_T(T)$ of ``exact order $N$'' on $C_T/T$ which generates the subgroup $G_T$ of $C_T$, i.e., such that we have an equality of Cartier divisors in $C_T$:
\begin{equation}\label{eq:1.4.1.2}
G_T = \sum_{a=1}^{N} [aP].
\end{equation}
Clearly if a point $P \in C(S)$ has ``exact order $N$'' in $C/S$, and generates a subgroup $G$ of rank $N$, then for any change of base $T \to S$ the induced point $P_T \in C(T) = C_T(T)$ has ``exact order $N$'' in $C_T/T$, and it generated the subgroup $G_T$ of $C_T/T$.

\begin{lemma}\label{lem:1.4.2}
If $P \in C(S)$ has ``exact order $N$,'' then $NP = 0$.
\end{lemma}

\begin{proof}
Any finite locally free commutative group-scheme of rank $N$ is known to be killed by $N$ (cf.~\cite{OortTate}). Therefore every section, in particular $P$, of the effective Cartier divisor $\sum_{1}^{N} [aP]$ is killed by $N$.
\end{proof}

\begin{caution}\label{caution:1.4.3}
If $S$ is an $F_p$-scheme, then for every integer $n \ge 1$ the zero-section $0 \in C(S)$ has ``exact order $p^n$''; indeed it ``generates'' the subgroup $\Ker(F^n : C \to C^{(p^n)})$ of rank $p^n$. This example shows that a given point can have many different ``exact orders.''
\end{caution}

\begin{lemma}\label{lem:1.4.4}
Suppose that $N$ is invertible on $S$. Let $P \in C(S)$ be a point killed by $N$. Then the following conditions are equivalent.
\begin{enumerate}[(1)]
\item $P$ has ``exact order $N$'' in $C/S$.
\item For every geometric point $\Spec(k) \to S$ of $S$, the induced point $P_k \in C(k) = C_k(k)$ has ``exact order $N$'' in $C_k/k$.
\item For every geometric point $\Spec(k) \to S$ of $S$, the induced point $P_k \in C(k)$ has exact order $N$ in the usual sense that $N$ is the least positive integer which kills $P_k$, i.e., the $N$ points $\{aP_k\}$, $a = 1, \ldots, N$ are all distinct in $C(k)$.
\item The effective Cartier divisor in $C/S$
\[
\sum_{a=1}^{N} [aP]
\]
in $C/S$ is finite etale over $S$.
\item The unique $S$-group homomorphism $\Z/\NZ \to C$ which maps ``$1$''$\in \Z/\NZ$ to $P \in C(S)$ defines a closed $S$-immersion
\[
\Z/\NZ \hookrightarrow C
\]
which identifies the constant $S$-scheme $\Z/\NZ$ with the Cartier divisor $\sum [aP]$:
\[
\Z/\NZ \xrightarrow{\ \sim\ } \sum_{a=1}^{N} [aP].
\]
\end{enumerate}
\end{lemma}

\begin{proof}
(1) $\Rightarrow$ (2), because the property of having ``exact order $N$'' is preserved under arbitrary changes of base $T \to S$.

(2) $\Rightarrow$ (3). Let $G$ be the rank $N$ subgroup-scheme of $C_k$ ``generated by'' $P_k$. Then $G$ is finite flat of rank $N$ over $k$. Because $N$ is invertible in $S$, it is invertible in $k$, and therefore $G$ is automatically finite etale over $k$ of rank $N$. Therefore as a Cartier divisor in $C_k$, $G$ consists of a uniquely determined set of $N$ distinct points. The equality of Cartier divisors in $C_k$
\[
G = \sum_{a=1}^{N} [aP_k]
\]
shows that the $N$ points $\{aP_k\}$, $a = 1, \ldots, N$, are all distinct in $C_k$, as required.

(3) $\Leftrightarrow$ (4). Let us denote by $D$ the effective Cartier divisor
\[
D = \sum_{a=1}^{N} [aP]
\]
in $C/S$. As a scheme over $S$, $D$ is finite and locally free of rank $N$ over $S$. It is finite etale over $S$ if and only if its discriminant (determinant of the matrix $\tr(e_i e_j)\cdots$) is invertible on $S$. This holds if and only if for all geometric points $\Spec(k) \to S$, the Cartier divisor in $C_k$
\[
D_k = \sum_{a=1}^{N} [aP_k]
\]
is finite etale over $k$, i.e., if and only if (3) holds.

(3) $\Leftrightarrow$ (5). The morphism of $S$-groups $\Z/\NZ \to C$ mapping $1$ to $P$ certainly factors through the effective Cartier divisor
\[
D = \sum_{a=1}^{N} [aP],
\]
thus defining an $S$-morphism
\[
(\Z/\NZ)_S \to D.
\]
Both source and target of this morphism are finite locally-free over $S$ of the same rank $N$, so locally on $S$, say $S = \Spec(R)$, the map on coordinate rings is given by an $N \times N$ matrix over $R$. Our map is an isomorphism if and only if this matrix has determinant a unit in $R$. Therefore our map is an isomorphism if and only if for all geometric points $\Spec(k) \to S$ of $S$, the $k$-morphism
\[
\Z/\NZ \xrightarrow{\ 1 \mapsto P_k\ } D_k = \sum_{a=1}^{N} [aP_k]
\]
is an isomorphism, i.e., if and only if the $N$ points $\{aP_k\}$, $a = 1, \cdots, N$, are all distinct.

(5) $\Rightarrow$ (1). Indeed if (5) holds, then $\sum_{a=1}^{N} [aP]$ is endowed with the structure of subgroup of $C/S$.
\end{proof}

\begin{remark}\label{rem:1.4.5}
If we do not assume $N$ invertible on $S$, we still have the implications
\[
(3) \Longleftrightarrow (4) \Longleftrightarrow (5) \Longrightarrow (1) \Longrightarrow (2).
\]
A point $P \in C(S)$ which is killed by $N$ and which satisfies any of the equivalent conditions (3), (4), (5), might be called an ``etale point of exact order $N$.''
\end{remark}

% ------------------------------------------------------------------------------
% Section 1.5
% ------------------------------------------------------------------------------
\section{A mild generalization: $A$-structures and $A$-generators}
\label{sec:1.5}

\kmitem\label{item:1.5.1}
As before, $C/S$ is a smooth commutative group-scheme over $S$ of relative dimension one (cf.~\ref{item:1.4.1}). Let $A$ be an ``abstract'' finite abelian group. A homomorphism of abstract groups
\begin{equation}\label{eq:1.5.1.1}
\phi : A \to C(S)
\end{equation}
is said to be an $A$-structure on $C/S$ if the effective Cartier divisor $D$ in $C/S$ of degree $= \#(A)$ defined by
\[
D = \sum_{a \in A} [\phi(a)]
\]
is a subgroup $G$ of $C/S$. We call this subgroup $G$ the $A$-subgroup of $C/S$ generated by $\phi$, and we call $\phi$ an $A$-generator of $G$. We say that a closed subgroup-scheme $G \subset C$ which is finite locally free over $S$ of rank $= \#(A)$ is an $A$-subgroup if, locally f.p.p.f.\ on $S$, it admits an $A$-generator, i.e., if there exists an f.p.l.p.f.\ $T \to S$ and an $A$-structure
\[
\phi : A \to C(T) = C_T(T)
\]
on $C_T$ which generates $G_T$.

Clearly if $\phi : A \to C(S)$ is an $A$-structure, generating $G$, then for any change of base $T \to S$, the composite
\[
\phi_T : A \to C(S) \to C(T) = C_T(T)
\]
is an $A$-structure on $C_T/T$, and it generates $G_T$.

\kmitem\label{item:1.5.2}
When we take $A$ to be $\Z/\NZ$, we recover the notion of a point of ``exact order $N$''; for a homomorphism
\[
\phi : \Z/\NZ \to C(S)
\]
is a $\Z/\NZ$-structure on $C/S$ if and only if the point $P = \phi(1) \in C(S)$ is a point of ``exact order $N$'' in $C(S)$. Similarly, a finite locally free of rank $N$ subgroup-scheme of $C/S$ is a $\Z/\NZ$-subgroup if and only if it is a cyclic subgroup.

The analogue of Lemma~\ref{lem:1.4.4} is the following, whose proof, entirely similar, is left to the reader.

\begin{lemma}\label{lem:1.5.3}
Let $C/S$ be as above, $A$ a finite abelian ``abstract'' group, and $\phi : A \to C(S)$ a group homomorphism. Consider the following conditions:
\begin{enumerate}[(1)]
\item $\phi$ is an $A$-structure on $C/S$.
\item For every geometric point $\Spec(k) \to S$ of $S$, the induced homomorphism $\phi_k : A \to C(k) = C_k(k)$ is an $A$-structure on $C_k/k$.
\item For every geometric point $\Spec(k) \to S$ of $S$, the induced homomorphism $\phi_k : A \to C(k)$ is injective.
\item The effective Cartier divisor in $C/S$
\[
\sum_{a \in A} [\phi(a)]
\]
is finite etale over $S$.
\item $\phi$ defines a closed $S$-immersion
\[
A \hookrightarrow C
\]
which identifies the constant $S$-scheme $A$ with the Cartier divisor $\sum [\phi(a)]$:
\[
A \xrightarrow{\ \sim\ } \sum_{a \in A} [\phi(a)].
\]
\end{enumerate}
We always have the implications
\[
(3) \Longleftrightarrow (4) \Longleftrightarrow (5) \Longrightarrow (1) \Longrightarrow (2),
\]
and, if $\#(A)$ is invertible on $S$, then all of (1), (2), (3), (4), (5) are equivalent.
\end{lemma}

% ------------------------------------------------------------------------------
% Section 1.6
% ------------------------------------------------------------------------------
\section{General representability theorems for \texorpdfstring{$A$}{A}-structures and \texorpdfstring{$A$}{A}-generators}\label{sec:1.6}

\kmitem\label{item:1.6.1}
As before, $C/S$ is a smooth commutative group-scheme over $S$ of relative dimension one (cf.~\ref{item:1.4.1}), and $A$ is a fixed ``abstract'' finite abelian group. For every integer $N \ge 1$, we denote by $C[N]$ the $S$-subgroup-scheme of $C/S$ of points of order $N$: for any $S$-scheme $T \to S$, we have
\[
C[N](T) = \Ker \text{ of } N : C(T) \to C(T).
\]
Let us fix an isomorphism of abelian groups
\begin{equation}\label{eq:1.6.1.1}
A \simeq \Z/N_1\Z \times \cdots \times \Z/N_r\Z.
\end{equation}
Then the functor on $S$-schemes $\Hom_{S\text{-}\mathrm{gp}}(A, C)$ defined by
\begin{equation}\label{eq:1.6.1.2}
T \mapsto \Hom_{\mathrm{gp}}(A, C(T)),
\end{equation}
is represented by the $S$-scheme of finite presentation
\begin{equation}\label{eq:1.6.1.3}
C[N_1] \underset{S}{\times} C[N_2] \underset{S}{\times} \cdots \underset{S}{\times} C[N_r].
\end{equation}

\begin{proposition}\label{prop:1.6.2}
Hypotheses and notations as above, the functor $A\text{-}\mathrm{Str}(C/S)$ on $S$-schemes
\[
T \mapsto \text{the set of $A$-structures on } C_T/T
\]
is represented by a closed subscheme of $\Hom_{S\text{-}\mathrm{gp}}(A, C) \simeq C[N_1] \underset{S}{\times} \cdots \underset{S}{\times} C[N_r]$ definable locally by $1 + \#(A) + (\#(A))^2$ equations.
\end{proposition}

\begin{proof}
By definition, the functor $A\text{-}\mathrm{Str}(C/S)$ is a sub-functor of the representable functor $\Hom_{S\text{-}\mathrm{gp}}(A, C)$. It is represented by the closed sub-scheme of $\Hom_{S\text{-}\mathrm{gp}}(A, C)$ over which the effective Cartier divisor of degree $\#(A)$ in $C \times_S \Hom_{S\text{-}\mathrm{gp}}(A, C)\allowbreak / \Hom_{S\text{-}\mathrm{gp}}(A, C)$
\[
\sum_{a \in A} [\phi_{\mathrm{univ}}(a)],
\]
attached to the universal homomorphism $\phi_{\mathrm{univ}} : A \to C$, is a subgroup. The asserted result now follows from \ref{cor:1.3.7}.
\end{proof}

\begin{corollary}\label{cor:1.6.3}
Let $N$ be the exponent of the group $A$. If the $S$-morphism
\[
[N] : C \to C
\]
is finite, then $A\text{-}\mathrm{Str}(C/S)$ is represented by a finite $S$-scheme of finite presentation.
\end{corollary}

\begin{proof}
Then $C[N_1] \underset{S}{\times} \cdots \underset{S}{\times} C[N_r]$ is itself a closed sub-scheme of finite presentation of the finite $S$-scheme $C[N] \underset{S}{\times} \cdots \underset{S}{\times} C[N]$.
\end{proof}

\begin{proposition}\label{prop:1.6.4}
Hypotheses and notations as above, suppose that the exponent $N$ of the group $A$ is invertible on $S$, [and that the $S$-group $C[N]$ is finite over $S$]. Then the $S$-scheme $A\text{-}\mathrm{Str}(C/S)$ is [finite] etale over $S$ (possibly empty!).
\end{proposition}

\begin{proof}
Because $N$ is invertible on $S$, and $C/S$ is smooth, the homomorphism
\[
[N] : C \to C
\]
is etale. Therefore $C[N]$ is etale over $S$, [hence finite etale over $S$]. The functor $A\text{-}\mathrm{Str}(C/S)$ is known to be [finite over $S$] of finite presentation, so it suffices to check that it is formally etale. Thus let $T$ be an $S$-scheme, $T_0 \subset T$ a closed subscheme defined by a nilpotent ideal, and
\[
\phi_0 : A \to C(T_0)
\]
an $A$-structure on $C_{T_0}/T_0$. We must show that $\phi_0$ extends uniquely to a homomorphism
\[
\phi : A \to C(T)
\]
which defines an $A$-structure on $C_T/T$. To see this, we argue as follows. Because $A$ has exponent $N$, $\phi_0$ automatically factors through $C[N]$:
\[
\phi_0 : A \to C[N](T_0).
\]
Because $C[N]$ is etale over $S$, $\phi_0$ extends uniquely to a homomorphism
\[
\phi : A \to C[N](T).
\]
To see that this $\phi$ defines an $A$-structure on $C_T/T$, it suffices, $N$ being invertible, to test after pulling back to each geometric point of $T$. But $T$ and $T_0$ have precisely the same geometric points, and at each of them $(\phi)_k$ is none other than $(\phi_0)_k$, which by hypothesis is an $A$-structure on $C_k/k$.
\end{proof}

\begin{proposition}\label{prop:1.6.5}
Hypotheses and notations as above, let $G \subset C$ be a closed $S$-subgroup-scheme of $C$ which is finite locally-free over $S$ of rank $= \#(A)$. The functor $A\text{-}\mathrm{Gen}(G/S)$ on $S$-schemes defined by
\[
T \mapsto \text{the set of $A$-generators of the subgroup $G_T$ in } C_T/T
\]
is represented by a finite $S$-scheme of finite presentation, which is a closed subscheme of the finite $S$-scheme of finite presentation $\Hom_{S\text{-}\mathrm{gp}}(A, G)$, defined locally by $\#(A)$ equations.
\end{proposition}

\begin{proof}
A group homomorphism $\phi : A \to C(S)$ is an $A$-generator of $G$ if and only if the following two conditions are satisfied:
\begin{enumerate}[(1)]
\item $\phi$ maps $A$ to the subgroup $G(S)$ of $C(S)$.
\item We have an equality of Cartier divisors $G = \sum_{a \in A} [\phi(a)]$.
\end{enumerate}
Therefore $A\text{-}\mathrm{Gen}(G/S)$ is represented by the closed subscheme of $\Hom_{S\text{-}\mathrm{gp}}(A, G)$ over which the effective Cartier divisor $D_{\mathrm{univ}}$ in $C \times_S \Hom_{S\text{-}\mathrm{gp}}(A, G)$ over $\Hom_{S\text{-}\mathrm{gp}}(A, G)$ associated to $\phi_{\mathrm{univ}} : A \to G$
\[
D_{\mathrm{univ}} = \sum_{a \in A} [\phi_{\mathrm{univ}}(a)]
\]
satisfies
\[
D_{\mathrm{univ}} = G.
\]
The $S$-scheme $\Hom_{S\text{-}\mathrm{gp}}(A, G)$ is itself a finite $S$-scheme of finite presentation, being none other than $G[N_1] \underset{S}{\times} \cdots \underset{S}{\times} G[N_r]$. The result now follows from \ref{cor:1.3.5}.
\end{proof}

\begin{remark}\label{rem:1.6.6}
The notation ``$A\text{-}\mathrm{Gen}(G/S)$'' makes no reference to the ambient $C/S$, although the definition of an $A$-generator for $G$ very much involves the ambient space, as the space in which the equality
\[
G = \sum_{a \in A} [\phi(a)]
\]
of Cartier divisors takes place. We will see later (Section~\ref{sec:1.9} ff.) that there is in fact an intrinsic notion of an $A$-generator for a finite locally-free commutative group-scheme $G/S$ which coincides with above notion whenever $G$ is embeddable as a subgroup of a smooth one-parameter group $C/S$.
\end{remark}

\begin{proposition}\label{prop:1.6.7}
Hypotheses and notations as in the preceding proposition, if $G$ is finite etale over $S$ (e.g., if its rank is invertible on $S$), then the functor $A\text{-}\mathrm{Gen}(G/S)$ is represented by a finite etale $S$-scheme.
\end{proposition}

\begin{proof}
Just as in the proof of \ref{prop:1.6.4}, it suffices to show that the functor $A\text{-}\mathrm{Gen}(G/S)$ is formally etale over $S$. But for $G$ etale over $S$, this is none other than the functor $\Isom_{S\text{-}\mathrm{gp}}(A, G)$, which for any two etale $S$-groups is itself formally etale over $S$.
\end{proof}

% ------------------------------------------------------------------------------
% Section 1.7
% ------------------------------------------------------------------------------
\section{Factorization into prime powers of $A$-structures and $A$-generators}\label{sec:1.7}

\kmitem\label{item:1.7.1}
Let $A_1$ and $A_2$ be finite abelian groups, of orders $N_1$ and $N_2$ respectively. Let $A = A_1 \times A_2$, and denote its order by $N = N_1 N_2$. Given any abelian group $B$, and any group homomorphism
\begin{equation}\label{eq:1.7.1.1}
\phi : A \to B
\end{equation}
there are unique group homomorphisms
\begin{equation}\label{eq:1.7.1.2}
\phi_i : A_i \to B, \quad i = 1, 2,
\end{equation}
such that
\begin{equation}\label{eq:1.7.1.3}
\phi(a_1, a_2) = \phi_1(a_1) + \phi_2(a_2).
\end{equation}

\begin{proposition}\label{prop:1.7.2}
Suppose that $N_1$ and $N_2$ are relatively prime. Let $C/S$ be a smooth commutative $S$-group of relative dimension one (cf.~\ref{item:1.4.1}). Then a homomorphism
\[
\phi : A \to C(S)
\]
is an $A$-structure on $C/S$ if and only if the two homomorphisms
\[
\phi_i : A_i \to C(S), \quad i = 1, 2,
\]
are each $A_i$-structures, $i = 1, 2$, on $C/S$. If this is the case, then the groups $G, G_1, G_2$ generated by $\phi, \phi_1, \phi_2$ respectively are related by $G = G_1 \underset{S}{\times} G_2$.
\end{proposition}

\begin{proof}
Suppose first that $\phi$ is an $A$-structure on $C/S$. Let $G \subset C$ be the $A$-subgroup generated by $\phi$. Then $G$ is finite locally-free over $S$ of rank $N = N_1 N_2$. Because $G$ is killed by $N$, and $N = N_1 N_2$ with $N_1$ and $N_2$ relatively prime, we have a canonical product decomposition (as abelian group valued functors on $S$-schemes)
\[
G \xleftarrow{\ \sim\ } G[N_1] \underset{S}{\times} G[N_2].
\]
Then $G[N_1]$ and $G[N_2]$ are finite locally-free over $S$ of ranks $N_1$ and $N_2$ respectively. To see this, notice first that $G[N_1]$ is representable by a closed subscheme of $G$ which is of finite presentation over $S$, simply because $G[N_1]$ is the kernel of an $S$-homomorphism $([N_1] : G \to G)$ between $S$-group-schemes of finite presentation. Because $G$ is finite over $S$, $G[N_1]$ is a finite $S$-scheme of finite presentation. Because $N = N_1 N_2$ with $N_1$ and $N_2$ relatively prime, $G[N_1]$ is a direct factor of $G$, corresponding to the $S$-projector on $G$
\[
P_1 + P_2 \mapsto P_1.
\]
Therefore the sheaf of $S$-algebras defining $G[N_1]$ is an $S$-direct factor of that defining $G$, so flat over $S$. Therefore $G[N_1]$ is finite flat over $S$ of finite presentation. By \cite[IV, 8.9.1 and 11.2.6.1]{EGA} we may reduce to the case $S$ noetherian to conclude that $G[N_1]$ is finite locally free over $S$. Similarly $G[N_2]$ is finite locally free over $S$. To show that $G[N_1]$ and $G[N_2]$ are locally free of the asserted ranks $N_1$ and $N_2$, respectively, it suffices to treat the case when $S$ is $\Spec(k)$ with $k$ an algebraically closed field. The product decomposition
\[
G[N_1] \underset{S}{\times} G[N_2] \xrightarrow{\ \sim\ } G
\]
shows that
\[
\rank(G[N_1]) \cdot \rank(G[N_2]) = \rank(G) = N_1 N_2.
\]
Because $G[N_i]$ is killed by $N_i$, the rank of $G[N_i]$ divides a power of $N_i$ (cf.~\cite[IV, \S3, 5.3--9]{DeGa}). Therefore the above product-formula for ranks forces $\rank(G[N_i]) = N_i$, as required.

Because $N_1$ and $N_2$ are relatively prime, we may, localizing on $S$ assume one of them, say $N_1$, is invertible on $S$. Then $G[N_1]$ is finite etale over $S$, of rank $N_1$.

We first show that $\phi_1 = \phi|_{A_1}$ is an $A_1$-structure on $C/S$, and that it $A_1$-generates $G[N_1]$. To show that $\phi_1$ is an $A_1$-structure, it suffices by \ref{lem:1.5.3} to show that for every geometric point $\Spec(k) \to S$ of $S$, the $N_1$ points $\{\phi(a_1)_k\}$, $a_1 \in A_1$, are all distinct in $G(k)$. To see this, notice that because $N_1$ is invertible in $k$ we have
\begin{align*}
G(k) &= G[N_1](k) \times G[N_2](k) \\
&= (\text{a group of order } N_1) \times (\text{a group killed by } N_2).
\end{align*}
Therefore $G(k)$ contains precisely $N_1$ distinct points killed by $N_1$. But the expression of $G$ as a Cartier divisor
\[
G = \sum_{a_1, a_2} [\phi(a_1) + \phi(a_2)]
\]
shows that the group $G(k)$ is exhausted by the points
\[
\{\phi(a_1)_k + \phi(a_2)_k\}, \quad a_1 \in A_1, \ a_2 \in A_2.
\]
But of these points, only those for which $\phi(a_2)_k = 0$ are killed by $N_1$. Therefore the $N_1$ points
\[
\{\phi(a_1)_k\}, \quad a_1 \in A_1
\]
must exhaust the point killed by $N_1$ in $G(k)$. Therefore these $N_1$ points are all distinct. This shows that $\phi_1$ is an $A_1$-structure on $C/S$. The map $\phi_1 : A_1 \to G$ clearly factors through $G[N_1]$, inducing
\[
\phi_1 : A_1 \to G[N_1].
\]
It remains to show that this last map is an isomorphism. As its source and target are finite flat of the same rank $N_1$ over $S$, it suffices to check that for every geometric point $\Spec(k) \to S$ of $S$, the induced map on $k$-valued points is bijective. But it is precisely this bijectivity
\[
\phi_1 : A_1 \xrightarrow{\ \sim\ } G(k)[N_1]
\]
which was established above.

We next must show that $\phi_2 = \phi|_{A_2}$ is an $A_2$-structure on $C/S$, and that it $A_2$-generates $G[N_2]$. We have already proved that $\phi_1$ defines an isomorphism
\[
\phi_1 : A_1 \xrightarrow{\ \sim\ } \sum_{a_1 \in A_1} [\phi_1(a_1)] = G[N_1].
\]
Let us denote by $D_2$ the effective Cartier divisor in $C/S$ defined by
\[
D_2 = \sum_{a_2 \in A_2} [\phi(a_2)].
\]
Then the decomposition
\[
G = \sum [\phi(a_1) + \phi(a_2)]
\]
may be rewritten
\[
G = \sum_{a_1 \in A_1} \mathrm{Trans}(\phi(a_1))^*(D_2),
\]
a decomposition in which the $N_1$ translates of $D_2$ which occur are everywhere disjoint (say as sub-functors of $C/S$ on $S$-schemes). Therefore as an $S$-scheme, $G$ is isomorphic to the disjoint union,
\[
G \simeq \coprod_{a_1 \in A_1} \mathrm{Trans}(\phi(a_1))^*(D_2).
\]
We must show that $G[N_2] = D_2$. Because both sides have the same degree, it suffices to show that $G[N_2] \le D_2$. For this, it suffices to show that for any point $P \in G[N_2](T)$, with $T$ any connected $S$-scheme, we have $P \in D_2(T)$. Certainly $P \in G(T)$, and, because $T$ is connected, the disjoint union expression of $G$ shows that
\[
G(T) = \coprod_{a_1 \in A_1} \mathrm{Trans}(\phi(a_1))^*(D_2(T)).
\]
In other words, there is an element $a_1 \in A_1$ such that
\[
P - \phi(a_1) = \text{an element of } D_2(T).
\]
We must show that $\phi(a_1) = 0$. For this, it suffices to make the change of base to any geometric point $\Spec(k) \to T$ of $T$. Then we find
\[
P_k - \phi(a_1)_k = \text{an element of } D_2(k) = \phi(a_2)_k \text{ for some } a_2 \in A_2.
\]
As both $P_k$ and $\phi(a_2)_k$ are killed by $N_2$, while $\phi(a_1)$ is killed by $N_1$, this is impossible unless $\phi(a_1) = 0$.

We now prove the converse. We are given
\[
\phi : A \to C(S)
\]
such that $\phi_i = \phi|_{A_i}$ is an $A_i$-structure, $i = 1, 2$, on $C/S$. We must show that $\phi$ is an $A$-structure. Because $N = N_1 N_2$ with $N_1$ and $N_2$ relatively prime, we may, by localizing on $S$, suppose that one of $N_1$ and $N_2$, say $N_1$, is invertible on $S$. Let $G_i \subset C$ be the $A_i$-subgroup generated by $\phi_i$, for $i = 1, 2$. Because $N_1$ is invertible on $S$, we know that $\phi_1$ defines an $S$-group isomorphism
\[
\phi_1 : A_1 \xrightarrow{\ \sim\ } G_1.
\]
Because $G_1$ and $G_2$ have relatively prime ranks, the sum map defines a closed immersion
\[
G_1 \underset{S}{\times} G_2 \hookrightarrow C,
\]
whose image $G$ is the disjoint union of the $N_1$ translates of $G_2$ by the given sections of $G_1$:
\[
G_1 \underset{S}{\times} G_2 \xrightarrow{\ \sim\ } G = \coprod_{a_1 \in A_1} \mathrm{Trans}(\phi(a_1))^*(G_2).
\]
Because these translates are everywhere disjoint in $C/S$, their union in $C$ is just their \emph{sum} as Cartier divisors, so that
\[
G_1 \underset{S}{\times} G_2 \xrightarrow{\ \sim\ } G = \sum_{a_1 \in A_1} \mathrm{Trans}(\phi(a_1))^*(G_2).
\]
By hypothesis on $\phi_2$ we have an equality of Cartier divisors in $C/S$,
\[
G_2 = \sum_{a_2 \in A_2} [\phi(a_2)],
\]
so that we find
\begin{align*}
G_1 \underset{S}{\times} G_2 \xrightarrow{\ \sim\ } G &= \sum_{a_1 \in A_1, a_2 \in A_2} [\phi(a_1) + \phi(a_2)] \\
&= \sum_{a \in A} [\phi(a)],
\end{align*}
as required.
\end{proof}

\begin{corollary}\label{cor:1.7.3}
Let $A$ be a finite abelian group, $N$ its order,
\[
N = \prod_{i=1}^r p_i^{n_i}
\]
the factorization of $N$ as a product of powers of distinct primes, and
\[
A = A_1 \times \cdots \times A_r
\]
the corresponding decomposition of $A$ into the product of its $p_i$-primary subgroups $A_i$. Let $C/S$ be a smooth commutative group-scheme of relative dimension one (cf.~\ref{item:1.4.1}). Then the decomposition $A = \prod A_i$ defines a canonical isomorphism of $S$-schemes
\[
A\text{-}\mathrm{Str}(C/S) \xrightarrow{\ \sim\ } A_1\text{-}\mathrm{Str}(C/S) \underset{S}{\times} \cdots \underset{S}{\times} A_r\text{-}\mathrm{Str}(C/S).
\]
For any closed $S$-subgroup-scheme $G \subset C/S$ which is finite and flat over $S$ of finite presentation and rank $N$, let
\[
G = G_1 \underset{S}{\times} \cdots \underset{S}{\times} G_r, \quad G_i \xlongequal{dfn} G[p_i^{n_i}]
\]
be the canonical factorization of $G$ into its $p$-primary components. Then we have a canonical isomorphism of $S$-schemes
\[
A\text{-}\mathrm{Gen}(G/S) \xrightarrow{\ \sim\ } A_1\text{-}\mathrm{Gen}(G_1/S) \underset{S}{\times} \cdots \underset{S}{\times} A_r\text{-}\mathrm{Gen}(G_r/S).
\]
\end{corollary}

% ------------------------------------------------------------------------------
% Section 1.8
% ------------------------------------------------------------------------------
\section{Full sets of sections}\label{sec:1.8}

\kmitem\label{item:1.8.1}
Let $S$ be a scheme, and $Z/S$ a finite flat $S$-scheme of finite presentation and rank $N \ge 1$. The standard reduction to the noetherian case (by \cite[IV, 8.9.1 and 11.2.6.1]{EGA}) shows that, equivalently, $Z/S$ is finite locally free over $S$ of rank $N \ge 1$. This means that for every affine $S$-scheme $\Spec(R) \to S$, the $R$-scheme $Z_R/R$ obtained from $Z/S$ by change of base is of the form $\Spec(B)$, where $B$ is an $R$-algebra which as an $R$-module is locally free of rank $N$. Any element $f \in B$ acts by multiplication to define an $R$-linear endomorphism of $B$. Because $B$ is a locally free $R$-module of rank $N$, we can speak of the characteristic polynomial of this endomorphism,
\begin{equation}\label{eq:1.8.1.1}
\det(T-f) = T^N - \trace(f) T^{N-1} + \cdots + (-1)^N \mathrm{Norm}(f),
\end{equation}
which is a monic polynomial in $R[T]$ of degree $N$.

\kmitem\label{item:1.8.2}
We say that a set of $N$ not-necessarily-distinct points $P_1, \dots, P_N$ in $Z(S)$ is a ``full set of sections of $Z/S$'' if either of the two following equivalent conditions is fulfilled:
\begin{enumerate}[(1)]
\item For every affine $S$-scheme $\Spec(R)$, and for every function $f$ on $Z_R$, i.e., for every $f \in B = H^0(Z_R, \cO)$, we have the equality of characteristic polynomials in $R[T]$
\[
\det(T-f) = \prod_{i=1}^N (T-f(P_i)).
\]
\item For every affine $S$-scheme $\Spec(R)$, and for every function $f$ on $Z_R$, we have the equality in $R$
\[
\mathrm{Norm}(f) = \prod_{i=1}^N f(P_i).
\]
\end{enumerate}
To see that (1) and (2) are equivalent, we argue as follows. We have the implication
\[
\text{(1) for } (R, f) \Longrightarrow \text{(2) for } (R, f)
\]
simply by putting $T = 0$ in (1), and we have the implication
\[
\text{(1) for } (R, f) \Longleftarrow \text{(2) for } (R[T], T-f)
\]
simply because the characteristic polynomial of $f \in B$ is just the norm of $T-f$ relative to $B \underset{R}{\otimes} R[T] / R[T]$.

\begin{lemma}\label{lem:1.8.3}
Let $Z/S$ be a finite etale $S$-scheme of rank $N$, and $P_1, \dots, P_N \in Z(S)$. Then the following conditions are equivalent:
\begin{enumerate}[(1)]
\item The $S$-morphism
\[
\underbrace{S \amalg S \amalg \cdots \amalg S}_{N\text{-fold disjoint union}} \longrightarrow Z
\]
defined by the $N$-sections $P_1, \dots, P_N$ is an isomorphism of $S$-schemes. If $S$ is connected, this is equivalent to saying that the set $Z(S)$ consists of $N$ elements, and that these elements are $P_1, \dots, P_N$.
\item $P_1, \dots, P_N$ form a full set of sections of $Z/S$.
\item For every geometric point of $S$, $\Spec(k) \to S$, the $N$ points $(P_i)_k \in Z(k)$ are all distinct.
\end{enumerate}
\end{lemma}

\begin{proof}
Trivially $(3) \Longleftrightarrow (1) \Longrightarrow (2)$. To prove $(2) \Longrightarrow (3)$, we may reduce to the case when $S$ is the spectrum of an algebraically closed field $k$. Then $Z$ is just $N$ distinct reduced points over $k$, say $Q_1, \dots, Q_N$. We must show that $P_1, \dots, P_N$ exhaust these $N$ points, i.e., that $P_1, \dots, P_N$ are all distinct. Let $f$ be any function on $Z$ which takes $N$ distinct values on $Q_1, \dots, Q_N$; by hypothesis, its characteristic polynomial is given by
\[
\prod (T-f(Q_i)) \xlongequal{dfn} \det(T-f) = \prod (T-f(P_i)).
\]
Therefore the $N$ values $f(P_i)$ are all distinct, and therefore the $P_i$ are all distinct.
\end{proof}

\begin{lemma}\label{lem:1.8.4}
Let $Z_1$ and $Z_2$ be two finite flat $S$-schemes of finite presentation, and ranks $N_1$ and $N_2$ respectively. Let
\[
\begin{cases}
P_1^{(1)}, \dots, P_{N_1}^{(1)} \in Z_1(S) \\
P_1^{(2)}, \dots, P_{N_2}^{(2)} \in Z_2(S)
\end{cases}
\]
be given sequences of $N_i$ points in $Z_i(S)$, for $i = 1, 2$. The following conditions are equivalent:
\begin{enumerate}[(1)]
\item For $i = 1, 2$, the $N_i$ $S$-valued points $P_1^{(i)}, \dots, P_{N_i}^{(i)}$ form a full set of sections of $Z_i/S$.
\item The $N_1 + N_2$ $S$-valued points $\{P_1^{(1)}, \dots, P_{N_1}^{(1)}, P_1^{(2)}, \dots, P_{N_2}^{(2)}\}$ of the disjoint union $Z_1 \amalg Z_2$ form a full set of sections of $(Z_1 \amalg Z_2)/S$.
\end{enumerate}
\end{lemma}

\begin{proof}
Locally on $S$, we have $S = \Spec(R)$, and $Z_i = \Spec(B_i)$ where $B_i$ is an $R$-algebra which as an $R$-module is free of rank $N_i$ for $i = 1, 2$. Then $Z_1 \amalg Z_2$ is $\Spec(B_1 \oplus B_2)$. For any element
\[
f = f_1 \oplus f_2 \quad \text{in} \quad B_1 \oplus B_2,
\]
we have
\[
\mathrm{Norm}(f) = \mathrm{Norm}(f_1) \mathrm{Norm}(f_2).
\]
For any of the $N_1 + N_2$ points $P_j^{(i)}$, we have
\[
f(P_j^{(i)}) = f_i(P_j^{(i)}).
\]
To show $(1) \Longrightarrow (2)$, write $f = f_1 \oplus f_2$ to compute $\mathrm{Norm}(f)$. To show $(2) \Longrightarrow (1)$, say for $Z_1$, notice that for $f_1 \in B_1$, we have
\[
\mathrm{Norm}(f_1) = \mathrm{Norm}(f_1 \oplus 1),
\]
and then use (2) to compute $\mathrm{Norm}(f_1 \oplus 1)$.
\end{proof}

\begin{lemma}\label{lem:1.8.5}
Let $Z_1/S$ and $Z_2/S$ be two\footnote{non-empty.} finite flat $S$-schemes of finite presentation, and ranks $N_1$ and $N_2$ respectively. Let
\begin{gather*}
P_1^{(1)}, \dots, P_{N_1}^{(1)} \in Z_1(S) \\
P_1^{(2)}, \dots, P_{N_2}^{(2)} \in Z_2(S)
\end{gather*}
be given sequences of $N_i$ points in $Z_i(S)$, for $i = 1, 2$. Suppose that $Z_1/S$ is finite etale. Then the following conditions are equivalent:
\begin{enumerate}[(1)]
\item For $i = 1, 2$, the $N_i$ points $P_1^{(i)}, \dots, P_{N_i}^{(i)}$ in $Z_i(S)$ form a full set of sections of $Z_i/S$.
\item The $N_1 \times N_2$ $S$-valued points $P_j^{(1)} \times P_k^{(2)}$ of $Z_1 \underset{S}{\times} Z_2$ form a full set of sections of $Z_1 \underset{S}{\times} Z_2$.
\end{enumerate}
\end{lemma}

\begin{proof}
To see that $(1) \Longrightarrow (2)$, notice that (1) for $Z_1$ insures that the $P_j^{(1)}$'s define an $S$-isomorphism
\[
S \amalg \cdots \amalg S \xrightarrow{\ \sim\ } Z_1,
\]
whence an $S$-isomorphism
\[
Z_2 \amalg \cdots \amalg Z_2 \xrightarrow{\ \sim\ } Z_1 \underset{S}{\times} Z_2.
\]
The result now follows from \ref{lem:1.8.4}.

It remains to prove $(2) \Longrightarrow (1)$. We first show that (1) holds for $Z_1/S$. For this, it suffices to treat the case when $S = \Spec(k)$ with $k$ algebraically closed. Consider the function $f_1 \otimes 1$ on $Z_1 \times_S Z_2$. Its characteristic polynomial is given by
\[
\begin{array}{c}
\det(T - f_1 \otimes 1) = (\det(T - f_1))^{N_2} \\
\rotatebox{90}{$=$} \\
\Biggl( \prod_{j=1}^{N_1} (T - f_1(P_j^{(1)})) \Biggr)^{\!N_2},
\end{array}
\]
the vertical equality in virtue of (2). Therefore we may infer
\[
\det(T - f_1) = \prod_{j=1}^{N_1} (T - f_1(P_j^{(1)})),
\]
since both sides are monic polynomials in $k[T]$. As this holds for any function $f_1$ on $Z_1$, it follows that the $P_j^{(1)}$'s are a full set of sections for $Z_1/k$ (compare the proof of \ref{lem:1.8.3}).

Once we know that the $P_j^{(1)}$'s are a full set of sections of $Z_1/S$, we have as before an $S$-isomorphism
\[
S \amalg \cdots \amalg S \xrightarrow{\ \sim\ } Z_1,
\]
whence
\[
Z_2 \amalg \cdots \amalg Z_2 \xrightarrow{\ \sim\ } Z_1 \underset{S}{\times} Z_2,
\]
defined by the sections $P_j^{(1)}$ of $Z_1/S$. It now follows from \ref{lem:1.8.4} that for each $1 \le j \le N_1$, the $N_2$ points $P_j^{(1)} \times P_k^{(2)}$, $k = 1, \dots, N_2$, are a full set of sections of $S \times_S Z_2 \xrightarrow{\sim} Z_2$, i.e., the points $P_k^{(2)}$, $k = 1, \dots, N_2$ are a full set of sections of $Z_2/S$.
\end{proof}

When we drop the hypothesis that $Z_1$ be etale over $S$, we have only (1) $\Longrightarrow$ (2). More precisely, we have:

\begin{lemma}\label{lem:1.8.6}
Let $Z_1/S$ and $Z_2/S$ be two finite flat $S$-schemes of finite presentation, and ranks $N_1$ and $N_2$ respectively. Suppose that for $i = 1, 2$, we are given sequences of $N_i$ $S$-valued points of $Z_i$:
\[
P_1^{(i)}, \dots, P_{N_i}^{(i)} \in Z_i(S) \quad \text{for } i = 1, 2,
\]
which are full sets of sections of $Z_i/S$ for $i = 1, 2$. Then the $N_1 \times N_2$ $S$-valued points $P_j^{(1)} \times P_k^{(2)}$ of $Z_1 \underset{S}{\times} Z_2$ form a full set of sections of $Z_1 \underset{S}{\times} Z_2/S$.
\end{lemma}

\begin{proof}
The points $P_j^{(1)}$ being a full set of sections of $Z_1/S$, they remain a full set of sections after any change of base $T \to S$ of the resultant situation $Z_1 \times_S T/T$. Taking for $T \to S$ the scheme $Z_2 \to S$, we see that the $P_j^{(1)}$'s give a full set of sections of $Z_1 \times_S Z_2 / Z_2$. The assertion now results from the transitivity of the norm in the tower
\[
\begin{tikzcd}
Z_1 \underset{S}{\times} Z_2 \arrow[d] \\
Z_2 \arrow[d] \\
S
\end{tikzcd}\qedhere
\]
\end{proof}

% ------------------------------------------------------------------------------
% Section 1.9
% ------------------------------------------------------------------------------
\section{Intrinsic $A$-structures and $A$-generators}\label{sec:1.9}

\begin{proposition}\label{prop:1.9.1}
Let $Z/S$ be a finite flat $S$-scheme of finite presentation and rank $N \ge 1$. Let $P_1, \dots, P_N \in Z(S)$ be a set of $N$ not-necessarily-distinct $S$-valued points of $Z$. There exists a unique closed subscheme $W$ of $S$ which is universal for the relation ``$P_1, \dots, P_N$ are a full set of sections of $Z/S$'' in the sense that for any $S$-scheme $T \to S$, the induced points $P_{1,T}, \dots, P_{N,T} \in Z(T)$ are a full set of sections of $Z_T/T$ if and only if the morphism $T \to S$ factors through $W$. Locally on $S$, $W$ is defined by finitely many equations.
\end{proposition}

\begin{proof}
The question is Zariski local on $S$, which we may therefore suppose affine, $S = \Spec(R)$. Further localizing on $S$, we may suppose that $Z = \Spec(B)$ where $B$ is an $R$-algebra which as an $R$-module is free of rank $N$. Let $b_1, \dots, b_N$ be an $R$-basis of $B$. Then the ``universal'' element $f \in B$ is the element
\[
f = \sum T_i b_i \in B \underset{R}{\otimes} R[T_1, \dots, T_N].
\]
The points $P_1, \dots, P_N$ form a full set of sections of $\Spec(B)/R$ if and only if this universal $f$ satisfies
\[
\mathrm{Norm}(f) = \prod f(P_i) \quad \text{in} \quad R[T_1, \dots, T_N],
\]
i.e., if and only if we have
\[
\mathrm{Norm}\Biggl( \sum T_i b_i \Biggr) = \prod_{j=1}^N \Biggl( \sum T_i b_i(P_j) \Biggr),
\]
equality inside $R[T_1, \dots, T_N]$. Both sides of this identity are homogeneous forms of degree $N$ in the variables $T_1, \dots, T_N$. The required closed sub-scheme $W \subset \Spec(R)$ is the one defined by the ideal in $R$ generated by the $\binom{2N-1}{N}$ coefficients of the difference of these two homogeneous forms.
\end{proof}

\begin{corollary}\label{cor:1.9.2}
Hypotheses and notations as above, suppose that $S$ is reduced. Then in order that $P_1, \dots, P_N$ form a full set of sections of $Z/S$, it is necessary and sufficient that for every geometric point $\Spec(k) \to S$ of $S$, the following condition be satisfied:
\par
For every function $f$ in the affine algebra $B \otimes k$ of $Z_k/k$, we have the equality in $k$
\[
\mathrm{Norm}(f) = \prod_{i=1}^N f((P_i)_k).
\]
\end{corollary}

\begin{proof}
The condition is clearly necessary. To see that it is sufficient, we may suppose $S = \Spec(R)$ with $R$ reduced, and $Z = \Spec(B)$ with $B$ an $R$-algebra which is a free $R$-module of rank $N$, say with basis $b_1, \dots, b_N$. We must show that
\[
\mathrm{Norm}\Biggl( \sum T_i b_i \Biggr) = \prod_{j=1}^N \Biggl( \sum T_i b_i(P_j) \Biggr)
\]
in the ring $R[T_1, \dots, T_N]$, which by hypothesis is reduced. To show this, it suffices to check equality of values at all geometric points of $\Spec(R[T_1, \dots, T_N])$, i.e., to check that for any homomorphism $R \to k$ and any $N$-tuple $(a_1, \dots, a_N) \in k^N$, the function $f = \sum a_i(b_i \otimes 1)$ satisfies the identity
\[
\mathrm{Norm}(f) = \prod_{i=1}^N f((P_i)_k). \qedhere
\]
\end{proof}

% ------------------------------------------------------------------------------
% Section 1.10
% ------------------------------------------------------------------------------
\section{Relation to Cartier divisors}\label{sec:1.10}

\begin{theorem}\label{thm:1.10.1}
Let $S$ be a scheme, $C/S$ a smooth curve (cf.~\ref{item:1.2.1}) and $Z \hookrightarrow C$ a closed subscheme which is finite flat over $S$ of finite presentation, and of rank $N \ge 1$. Let $P_1, \dots, P_N$ be a set of $N$ not-necessarily distinct points of $C(S)$. Then the following conditions are equivalent:
\begin{enumerate}[(1)]
\item We have an equality of effective Cartier divisors in $C/S$
\[
Z = \sum_{i=1}^N [P_i].
\]
\item The points $P_1, \dots, P_N$ all lie in the subset $Z(S)$ of $C(S)$, and they are a ``full set of sections'' of $Z/S$.
\end{enumerate}
\end{theorem}

\begin{proof}
A standard reduction reduces us to treating the case when $S = \Spec(R)$ with $R$ an artin local ring with algebraically closed residue field $k$. Let $z_1, \dots, z_r$ denote the distinct points of $Z(k)$. Then we have a canonical decomposition of $Z$ into a disjoint union
\[
Z = \coprod_{j=1,\dots,r} Z_j
\]
where $Z_j$ is the completion of $Z$ along $z_j$. Because $\Spec(R)$ is connected, we have a disjoint union decomposition
\[
Z(R) = \coprod_{j=1,\dots,r} Z_j(R).
\]
Under either of the conditions (1), (2) to be proven equivalent, the points $P_1, \dots, P_N$ all lie in $Z(R)$. Therefore we may partition them according to which subset $Z_j(R)$ of $Z(R)$ they lie in, say
\[
P_1^{(j)}, \dots, P_{n(j)}^{(j)} \quad \text{lie in} \quad Z_j, \quad j = 1, \dots, r.
\]
One verifies easily that in this situation of disjointness, each of the conditions (1), (2) holds for $Z$ and $P_1, \dots, P_N$ if and only if it holds for the $r$ situations ($Z_j$ and $P_1^{(j)}, \dots, P_{n(j)}^{(j)}$), $j = 1, \dots, r$.

Therefore it suffices to treat the case where $Z$ is itself local, concentrated at $z = z_1 \in Z(k) \subset C(k)$. Let $\widehat{\cO}_{C,z}$ denote the complete local ring of $C$ at $z$. In terms of a choice of uniformizing parameter $X$ at $z$, we obtain a noncanonical isomorphism
\[
\widehat{\cO}_{C,z} \xrightarrow{\ \sim\ } R[[X]].
\]
By Weierstrass preparation the effective Cartier divisor $Z$ of degree $N$ is defined inside $\Spec(R[[X]])$ by a unique monic equation $F(X)$ of degree $N$, all of whose lower coefficients lie in the maximal ideal of $R$
\begin{align*}
Z &\simeq \Spec(R[[X]]/(F(X))) \\
&\simeq \Spec(R[X]/(F(X)))
\end{align*}
\[
F(X) = X^N + (\text{lower terms with coefficients in } \max(R)).
\]
The effective Cartier divisor in $C/R$
\[
\sum [P_i]
\]
is also concentrated at $z$; it is defined inside $\Spec(R[[X]])$ by the monic polynomial $G(X)$ of degree $N$ defined by
\[
G(X) = \prod_{i=1}^N (X - X(P_i)).
\]
All of the lower coefficients of $G$ lie in $\max(R)$, because each $X(P_i) \in R$ is a root of $F$ and so must lie in $\max(R)$.

Because the effective Cartier divisors $Z$ and $\sum [P_i]$ are both supported at $z \in C$, they are equal in $C/R$ if and only if they are equal inside $\Spec(R[[X]])$. Thus (1) is equivalent to the equality of monic polynomials
\[
F(X) = \prod_{i=1}^N (X - X(P_i)).
\]
Because $Z$ is the finite free $R$-scheme
\[
Z = \Spec(R[X]/(F(X))),
\]
we are reduced to treating the case $S = \Spec(R)$, $C = \Spec(R[X])$, and $Z$ is defined by a monic polynomial $F(X)$ of degree $N$. Let $a_i = X(P_i)$ for $i = 1, \dots, N$. The assertion amounts to the following.
\end{proof}

\begin{lemma}\label{lem:1.10.2}
Let $R$ be a ring, $F(X) \in R[X]$ a monic polynomial of degree $N \ge 1$, and $a_1, \dots, a_N$ elements of $R$. Let $B$ denote the $R$-algebra $R[X]/(F(X))$. Then the following two conditions are equivalent:
\begin{enumerate}[(1)]
\item We have the factorization
\[
F(X) = \prod_{i=1}^N (X - a_i).
\]
\item For every $f \in B$, we have the factorization
\[
\det(T - f) = \prod_{i=1}^N (T - f(a_i)).
\]
\end{enumerate}
\end{lemma}

\begin{proof}
To see that $(2) \Longrightarrow (1)$, notice that for $R[X]/(F(X))$ we have the identity
\[
\det(T - X) = F(T).
\]
Therefore if (2) holds for $f = X$, we obtain
\[
\prod_{i=1}^N (T - a_i) = F(T),
\]
as required. To prove $(1) \Longrightarrow (2)$, we may reduce to the universal case
\begin{gather*}
R = \Z[A_1, \dots, A_N, B_0, \dots, B_{N-1}] \\
(a_1, \dots, a_N) = (A_1, \dots, A_N) \\
F(X) = \prod_{i=1}^N (X - A_i) \\
f = \sum_{i=0}^{N-1} B_i X^i,
\end{gather*}
in which the $A$'s, $B$'s are independent variables. As it is enough to verify that (2) holds after any injective extension of scalars $R \hookrightarrow R'$, we may pass to the fraction field of $R$ and thus reduce to the case in which $R$ is a field and $a_1, \dots, a_N$ are all distinct. Then we have an algebra isomorphism
\[
R[X]/\Biggl(\prod_{i=1}^N (X - a_i)\Biggr) \xrightarrow{\ \sim\ } \prod_{i=1}^N R[X]/(X - a_i),
\]
which makes $(1) \Longrightarrow (2)$ obvious.
\end{proof}

\begin{corollary}\label{cor:1.10.3}
Let $S$ be a scheme, $Z/S$ a finite flat $S$-scheme of finite presentation, and $P_1, \dots, P_r \in Z(S)$. Suppose that $Z/S$ is embeddable as a closed subscheme of a smooth curve $C/S$ (cf.~\ref{item:1.2.1}). Then there is at most one closed subscheme $W \hookrightarrow Z$ such that both of the following conditions hold:
\begin{enumerate}[(a)]
\item $W/S$ is locally free over $S$ of rank $r$.
\item $P_1, \dots, P_r \in W(S)$, and they form a ``full set of sections'' of $W/S$.
\end{enumerate}
This $W$ exists if and only if, locally f.p.p.f.\ on $S$, $P_1, \dots, P_r$ can be completed to a full set of sections $P_1, \dots, P_r, P'_{r+1}, \dots, P'_N$ of $Z/S$. Furthermore, there is a closed subscheme of $S$, defined locally by finitely many equations, which is universal for the existence of this $W$.
\end{corollary}

\begin{proof}
Inside $C/S$, form $W' = \sum [P_i]$. Then if $W$ exists, it is none other than $W'$, and the closed subscheme of $S$ over which $W$ exists is precisely the locus $W' \le Z$. The condition $W' \le Z$ inside $C$ is equivalent to the possibility, locally f.p.p.f.\ on $S$, of extending $P_1, \dots, P_r$ to a full set $P_1, \dots, P_r, P'_{r+1}, \dots, P'_N$ of sections of $Z/S$.
\end{proof}

\begin{remark}\label{rem:1.10.4}
This ``uniqueness of $W$'' is false without the ``embeddable in a curve'' hypothesis. For example, take $Z$ to be $\alpha_p \times \alpha_p$. Then the $p$ points $0, \dots, 0$, form a full set of sections of any of the $\PP^1$ of $\alpha_p$'s sitting inside $\alpha_p \times \alpha_p$.
\end{remark}

\kmitem\label{km:1.10.5}
Now let $G/S$ be a finite flat $S$-group-scheme of finite presentation, and rank $N \ge 1$. Let $A$ be a finite abelian ``abstract'' group of order $N$. We say that a group homomorphism
\begin{equation}\label{eq:1.10.5.1}
\phi \colon A \longrightarrow G(S)
\end{equation}
is an \emph{$A$-generator} of $G/S$ if the $N$ points $\{\phi(a)\}$, $a \in A$ are a ``full set of sections'' of $G/S$.

It results immediately from the preceding corollary that this notion of $A$-generator is compatible with the previous one. Explicitly, we have

\begin{proposition}\label{prop:1.10.6}
Let $C/S$ be a smooth commutative group-scheme over $S$ of relative dimension one (cf.~\ref{item:1.4.1}). Let $G \subset C$ be a closed $S$-subgroup-scheme which is finite flat over $S$ of rank $N$, and of finite presentation. Let $A$ be a finite ``abstract'' abelian group of order $N$, and let
\[
\phi \colon A \longrightarrow C(S)
\]
be a group homomorphism. Then the following conditions are equivalent:
\begin{enumerate}[(1)]
\item $\phi$ is an $A$-structure on $C/S$, and it $A$-generates $G/S$ in the sense of \ref{item:1.5.1}.
\item $\phi$ maps $A$ to the subgroup $G(S) \subset C[N](S)$, and the homomorphism $\phi \colon A \to G(S)$ is an $A$-generator of $G/S$ in the sense of \ref{km:1.10.5}.
\end{enumerate}
\end{proposition}

As an immediate corollary, we obtain

\begin{corollary}\label{cor:1.10.7}
Let $C/S$ be a smooth commutative group-scheme over $S$ of relative dimension one, and of finite presentation. Let $N \ge 1$ be an integer, and let $A$ be a finite abelian ``abstract'' group of order $N$. Then there is a natural functorial bijection
\[
\begin{array}{c}
\{A\text{-structures on } C/S\} \\
\big\downarrow \\
\{A\text{-structures on } C[N]/S\},
\end{array}
\]
which we define to mean pairs $(G, \phi)$ consisting of
\begin{enumerate}[(1)]
\item a closed $S$-subgroup-scheme $G \subset C[N]$ which is finite flat over $S$ of finite presentation and of rank $= N$
\item an $A$-generator $\phi$ of $G/S$ in the sense of \ref{km:1.10.5}.
\end{enumerate}
In particular, the notion of $A$-structure on $C/S$ may be described entirely in terms of the $S$-group-scheme $C[N]$.
\end{corollary}

\begin{corollary}\label{cor:1.10.8}
Suppose that $C/S$ and $C'/S$ are two smooth commutative group-schemes over $S$ of relative dimension one (cf.~\ref{item:1.4.1}). Suppose that for some integer $N \ge 1$, we have an isomorphism of $S$-groups
\[
C[N] \xrightarrow{\ \sim\ } C'[N].
\]
Then via \ref{cor:1.10.7} above we have an induced isomorphism of $S$-schemes
\[
A\text{\rm -Str}(C/S) \xrightarrow{\ \sim\ } A\text{\rm -Str}(C'/S)
\]
for any finite abelian ``abstract'' group $A$ of order $N$.
\end{corollary}

\begin{proof}
This results immediately from the description of $A$-structures in terms of data involving only the $N$-division points.
\end{proof}

\begin{corollary}\label{cor:1.10.9}
With the hypotheses and notations of \ref{cor:1.10.7} above, suppose in addition that the group-scheme $C[N]$ is finite and flat over $S$. Let
\[
\phi \colon A \longrightarrow C[N](S)
\]
be a group homomorphism. Then
\begin{enumerate}[(1)]
\item If there exists a closed $S$-subgroup-scheme $G \subset C[N]$ such that $(G, \phi)$ is an $A$-structure on $C[N]/S$, then $G$ is unique.
\item There exists a closed subscheme of $S$, defined locally by finitely many equations, which is universal for the existence of this $G$.
\end{enumerate}
\end{corollary}

\begin{proof}
Inside $C/S$, consider the Cartier divisor $W = \sum [\phi(a)]$, the sum over the $N$ elements $a \in A$. Then as Cartier divisors in $C/S$, we must have $G = W$. This gives uniqueness of $G$, and $G$ exists if and only if we first have $W \le C[N]$ and if, secondly, this $W$ inside $C[N]$ is a subgroup-scheme. As these conditions are successively defined, locally on $S$, by finitely many equations, we get the second assertion as well.
\end{proof}

\begin{remark}\label{rem:1.10.10}
Analogously, given an abstract finite abelian group $A$ of order $N$ and a finite flat commutative group-scheme of finite presentation $G$ over an arbitrary $S$, we can define an $A$-structure on $G/S$ to be a pair $(K, \phi)$ consisting of
\begin{itemize}
\item a finite flat $S$-subgroup-scheme $K \subset G$ which is locally free of rank $N = \# A$.
\item a homomorphism $\phi \colon A \to G(S)$ which lands in $K(S)$ and which is an ``$A$-generator'' of $K$.
\end{itemize}
The analogous result (with the same proof) is:
\end{remark}

\begin{lemma}\label{lem:1.10.11}
If $G/S$ is embeddable as a closed subscheme of a smooth curve $C/S$ (cf.~\ref{item:1.2.1}), then given a homomorphism
\[
\phi \colon A \longrightarrow G(S),
\]
there is at most one $K$ such that $(K, \phi)$ is an $A$-structure on $G/S$, and there is a closed subscheme of $S$, defined locally by finitely many equations, which is universal for the existence of this $K$. In particular, the functor $A\text{\rm -Str}(G/S)$ is represented by a closed subscheme of the finite $S$-scheme $\Hom(A, G)$.
\end{lemma}

\begin{proposition}\label{prop:1.10.12}
Let $G/S$ be a finite flat commutative group-scheme of finite presentation, and rank $N$. Let $A$ be a finite abelian group of order $N$, and
\[
\phi \colon A \longrightarrow G(S)
\]
a group homomorphism. Then the following conditions are equivalent:
\begin{enumerate}[(1)]
\item $\phi$ defines an isomorphism of $S$-group-schemes
\[
(A)_S \xrightarrow{\ \sim\ } G.
\]
\item $G/S$ is etale, and $\phi$ is an $A$-generator of $G/S$.
\end{enumerate}
If $S$ is connected, then (1) and (2) are each equivalent to
\begin{enumerate}[(3)]
\item $G/S$ is a constant group-scheme, and $\phi \colon A \to G(S)$ is an isomorphism of abstract groups.
\end{enumerate}
\end{proposition}

\begin{proof}
This is immediate from its set-theoretic analogue \ref{lem:1.8.3}.
\end{proof}

\begin{proposition}\label{prop:1.10.13}
Let $G/S$ be a finite flat commutative group-scheme of finite presentation, and rank $N$. Let $A$ be a finite abelian group of order $N$.
\begin{enumerate}[(1)]
\item Then the functor $A\text{\rm -Gen}(G/S)$ on $S$-schemes defined by
\[
T \longmapsto \left\{ \parbox{0.75\linewidth}{\centering the set of group homomorphisms $\phi \colon A \to G(T)$ such that the elements $\{\phi(a)\}$, $a \in A$ are a full set of sections of $G_T/T$} \right\}
\]
is representable by a finite $S$-scheme of finite presentation, namely by the closed subscheme of $\Hom_{S\text{\rm -gp}}(A, G)$ over which the image sections $\{\phi_{\mathrm{univ}}(a)\}$, $a \in A$ of the universal homomorphism $A \to G$ form a ``full set of sections.''
\item If $G/S$ is finite etale over $S$ of rank $N$, then by \ref{prop:1.10.12}, we have
\[
A\text{\rm -Gen}(G/S) \simeq \Isom_{S\text{\rm -gp}}((A)_S, G),
\]
so that over each connected component of $S$, $A\text{\rm -Gen}(G/S)$ is either empty or is a finite etale $\Aut(A)$-torsor.
\end{enumerate}
\end{proposition}

\begin{proposition}\label{prop:1.10.14}
Let $G/S$ be a finite flat commutative group-scheme of finite presentation, and rank $N$. Let $A$ be a finite abelian group of order $N$, and
\[
\phi \colon A \longrightarrow G(S)
\]
a group-homomorphism. Let
\[
N = p_1^{n_1} \cdots p_r^{n_r}
\]
be the factorization of $N$ into distinct prime powers, and let
\[
\begin{array}{r@{\ }l@{\qquad}r@{\ }l}
A &= \prod A_i & A_i &\xlongequal{dfn} A[p_i^{n_i}] \\[4pt]
G &= G_1 \underset{S}{\times} \cdots \underset{S}{\times} G_r & G_i &\xlongequal{dfn} G[p_i^{n_i}]
\end{array}
\]
be the corresponding decompositions of $A$ and $G$ into $p$-primary components. Then the following conditions are equivalent:
\begin{enumerate}[(1)]
\item $\phi$ is an $A$-generator of $G/S$;
\item for each $i = 1, \dots, r$, the induced homomorphism is an $A_i$-generator of $G_i/S$.
\end{enumerate}
\end{proposition}

\begin{proof}
If the number $r$ of distinct primes dividing $N$ is one, there is nothing to prove. If $r \ge 2$, then locally on $S$ all save at most one $p_i$, say $p_r$, is invertible. Then the groups $G_1, \dots, G_{r-1}$ are all finite etale, and the result follows from its set-theoretic analogue \ref{lem:1.8.5} applied successively to $G_1 \times (G_2 \times \cdots \times G_r)$, $G_2 \times \cdots \times G_r, \dots, G_{r-1} \times G_r$.
\end{proof}

\begin{corollary}\label{cor:1.10.15}
Hypotheses and notations as above, we have a canonical isomorphism of $S$-schemes
\[
A\text{\rm -Gen}(G/S) \xrightarrow{\ \sim\ } A_1\text{\rm -Gen}(G_1/S) \underset{S}{\times} \cdots \underset{S}{\times} A_r\text{\rm -Gen}(G_r/S).
\]
\end{corollary}

% ------------------------------------------------------------------------------
% Section 1.11
% ------------------------------------------------------------------------------
\section{Extensions of an etale group}\label{sec:1.11}

\kmitem\label{km:1.11.1}
Let $S$ be a connected scheme, and suppose given a short exact sequence (as f.p.p.f.\ abelian sheaves on the category of $S$-schemes, (cf.~\cite[III, Exp.~IV]{SGA}))
\begin{equation}\label{eq:1.11.1.1}
0 \longrightarrow H \longrightarrow G \longrightarrow E \longrightarrow 0
\end{equation}
of finite flat commutative $S$-group-schemes of finite presentation, with $E/S$ finite etale. Let $A$ be a finite abelian group whose order is the rank $G/S$, and let
\begin{equation}\label{eq:1.11.1.2}
\phi \colon A \longrightarrow G(S)
\end{equation}
be a group homomorphism. We will analyze exactly what it means, in terms of the given structure of extension on $G/S$, for $\phi$ to be a generator of $G/S$.

Consider first the composite homomorphism
\begin{equation}\label{eq:1.11.1.3}
A \longrightarrow G(S) \longrightarrow E(S),
\end{equation}
which we may view as a morphism between finite etale $S$-group-schemes
\begin{equation}\label{eq:1.11.1.4}
(A)_S \longrightarrow E.
\end{equation}
Such a morphism has as kernel a finite etale $S$-group, which is in fact constant (because it is a subgroup of the constant group $(A)_S/S$, and $S$ is connected), thus corresponding to an ``abstract'' subgroup
\begin{equation}\label{eq:1.11.1.5}
K \subset A.
\end{equation}
For any non-empty $S$-scheme $T \to S$, we have
\begin{equation}\label{eq:1.11.1.6}
K = \Ker(A \to E(T)).
\end{equation}
Thus we have a short exact sequence of abelian groups
\begin{equation}\label{eq:1.11.1.7}
0 \longrightarrow K \longrightarrow A \longrightarrow A/K \longrightarrow 0,
\end{equation}
and a commutative diagram with exact rows
\begin{equation}\label{eq:1.11.1.8}
\begin{tikzcd}[column sep=scriptsize]
0 \arrow[r] & K \arrow[r] \arrow[d, "\phi|K"'] & A \arrow[r] \arrow[d, "\phi"'] & A/K \arrow[r] \arrow[d, "\phi|(A/K)"'] & 0 \\
0 \arrow[r] & H(S) \arrow[r] & G(S) \arrow[r] & E(S)
\end{tikzcd}
\end{equation}
In terms of this commutative diagram, we have the following criterion.

\begin{proposition}\label{prop:1.11.2}
Hypotheses and notations as above, $\phi$ is an $A$-generator of $G/S$ if and only if both of the following conditions are satisfied:
\begin{enumerate}[(1)]
\item The order of $K$ is equal to the rank of $H/S$, and the induced homomorphism $\phi|K$ is a $K$-generator of $H/S$.
\item The order of $A/K$ is equal to the rank of $E/S$, and the induced homomorphism $\phi|(A/K)$ is an $A/K$-generator of $E/S$.
\end{enumerate}
\end{proposition}

\begin{proof}
Suppose first that $\phi$ is an $A$-generator of $G/S$. Then for every geometric point $\Spec(k) \to S$ of $S$, $\phi_k$ is an $A$-generator of $G_k/k$. Therefore for every such geometric point the homomorphism on $k$-valued points
\[
\phi_k \colon A \longrightarrow G(k)
\]
must be surjective. Over an algebraically closed field, the extension structure on $G$ provides a short exact sequence of groups of $k$-valued points
\[
0 \longrightarrow H(k) \longrightarrow G(k) \longrightarrow E(k) \longrightarrow 0,
\]
which sits in the commutative diagram with exact rows
\[
\begin{tikzcd}[column sep=small]
0 \arrow[r] & K \arrow[r] \arrow[d] & A \arrow[r] \arrow[d] & A/K \arrow[r] \arrow[d, hook] & 0 \\
0 \arrow[r] & H(k) \arrow[r] & G(k) \arrow[r] & E(k) \arrow[r] & 0.
\end{tikzcd}
\]
By the snake lemma, the right-most vertical arrow is an isomorphism
\[
A/K \xrightarrow{\ \sim\ } E(k).
\]
This shows that $A/K$ has order equal to the rank of $E/S$, and that $\phi|(A/K)$ defines an $A/K$-structure on $E/S$. Therefore (2) is proven.

If (2) holds, then any choice of \emph{set-theoretic} section
\[
\begin{tikzcd}
0 \arrow[r] & K \arrow[r] & A \arrow[r] & A/K \arrow[r] \arrow[l, bend right=40, "s"'] & 0
\end{tikzcd}
\]
defines a scheme-theoretic isomorphism of $S$-schemes
\[
H \underset{S}{\times} E \xrightarrow{\ \sim\ } G.
\]
It now results immediately from \ref{lem:1.8.5} that, given (2), then (1) holds if and only if $\phi$ is an $A$-structure on $G/S$.
\end{proof}

If we consider arbitrary extensions, we have only the following weaker result.

\begin{proposition}\label{prop:1.11.3}
Over an arbitrary scheme $S$, suppose given a short exact sequence (as f.p.p.f.\ abelian sheaves)
\begin{equation}\label{eq:1.11.3.1}
0 \longrightarrow G_1 \longrightarrow G \longrightarrow G_2 \longrightarrow 0
\end{equation}
of finite flat commutative $S$-group-schemes of finite presentation, and ranks $N_1$, $N$, and $N_2$ respectively.

Let $A_1$, $A$, and $A_2$ be finite abelian groups of orders $N_1$, $N$, and $N_2$ respectively, and suppose given a short exact sequence
\begin{equation}\label{eq:1.11.3.2}
0 \longrightarrow A_1 \longrightarrow A \longrightarrow A_2 \longrightarrow 0,
\end{equation}
and a commutative diagram of homomorphisms $\phi_1, \phi, \phi_2$:
\begin{equation}\label{eq:1.11.3.3}
\begin{tikzcd}[column sep=small]
0 \arrow[r] & A_1 \arrow[r] \arrow[d, "\phi_1"'] & A \arrow[r] \arrow[d, "\phi"'] & A_2 \arrow[r] \arrow[d, "\phi_2"'] & 0 \\
0 \arrow[r] & G_1(S) \arrow[r] & G(S) \arrow[r] & G_2(S)
\end{tikzcd}
\end{equation}
If for $i = 1, 2$ the homomorphisms $\phi_i$ are $A_i$-generators of $G_i/S$, then $\phi$ is an $A$-generator of $G/S$.
\end{proposition}

\begin{proof}
This is an immediate consequence of the transitivity of the norm for the finite locally free morphisms $G \to G_2$ and $G_2 \to S$, together with the observation that the fiber of $G \to G_2$ over any $S$-valued point $\phi_2(a_2)$ of $G_2$ is $G_1 \to S$. Therefore for $f$ any function on $G$, we have
\begin{align*}
N_{G/S}(f) &= N_{G_2/S}(N_{G/G_2}(f)) \\
&= \prod_{a_2 \in A_2} \{ (N_{G/G_2}(f))(\phi_2(a_2)) \} \\
&= \prod_{a_2 \in A_2} \biggl\{ \prod_{a \mapsto a_2} f(\phi(a)) \biggr\}. \qedhere
\end{align*}
\end{proof}

\begin{remark}\label{rem:1.11.4}
Here is an example to show that the converse of the preceding proposition is false. Let $S = \Spec(R)$ with $R$ an $\F_p$-algebra. For any $p$'th root of unity $\zeta \in \mu_p(R)$, we have a commutative diagram
\[
\begin{tikzcd}[column sep=small]
0 \arrow[r] & \Z/p\Z \arrow[r] \arrow[d, "\phi_1"'] & \Z/p\Z \oplus \Z/p\Z \arrow[r] \arrow[d, "\phi"'] & \Z/p\Z \arrow[r] \arrow[d, "\phi_2"'] & 0 \\
0 \arrow[r] & \mu_p(R) \arrow[r] & \mu_{p^2}(R) \arrow[r, "\times p"] & \mu_p(R)
\end{tikzcd}
\]
with
\[
\phi_1(a) = \zeta^a, \quad \phi(a, b) = \zeta^a, \quad \phi_2(b) = 1.
\]
Then $\phi$ is a $(\Z/p\Z \times \Z/p\Z)$-generator of $\mu_{p^2}/R$, because we have the polynomial identity
\begin{align*}
\prod_{a, b \in \Z/p\Z} (T - \phi(a, b)) &= \prod_a (T - \zeta^a)^p \\
&= \prod_a (T^p - 1) \\
&= T^{p^2} - 1.
\end{align*}
However, in order for $\phi_1$ to be a $\Z/p\Z$-structure on $\mu_p/R$, we must have the polynomial identity
\[
T^p - 1 = \prod_a (T - \zeta^a);
\]
in particular, $\zeta$ must be a root of the polynomial
\[
\Phi_p(X) = 1 + X + \dots + X^{p-1}
\]
as one sees by comparing coefficients of $T^{p-1}$. Thus if we take
\[
R = \F_p[X]/(X^p - 1), \quad \zeta = X,
\]
then $\phi_1$ is not a $\Z/p\Z$-generator of $\mu_p$.

A second example is provided by the product situation
\[
A = \Z/p\Z \oplus \Z/p\Z \xrightarrow{\ \phi\ } \alpha_p \times \alpha_p
\]
\[
\phi(a, b) = (0, bY)
\]
over an $\F_p$-algebra $R$, with $Y \in R$ satisfying $Y^p = 0$, $Y^{p-1} \neq 0$. Then $\phi$ and $\phi_1$ are generators, but not $\phi_2$.
\end{remark}

\kmitem\label{km:1.11.5}
Here is a final example in which $\phi$ is a generator, but neither $\phi_1$ nor $\phi_2$ is a generator. One takes the product situation
\[
A = \Z/p\Z \oplus \Z/p\Z \xrightarrow{\ \phi=0\ } \mu_p \times \mu_p,
\]
over the ring $\Z/p^2\Z$. In this situation, both $\phi_1$ and $\phi_2$ are the zero-map $\Z/p\Z \to \mu_p$, and this map is not a generator of $\mu_p$ over $\Z/p^2\Z$, simply because we do not have the required equality of Cartier divisors in $\mathbb{G}_m$ over $\Z/p^2\Z$:
\[
(X - 1)^p \neq (X^p - 1) \quad \text{in } (\Z/p^2\Z)[X].
\]
To see that $\phi$ is a generator of $\mu_p \times \mu_p$ over $\Z/p^2\Z$, we will verify the norm criterion: for any algebra $R$, and for any
\[
f(X, Y) \in R[X, Y]/(X^p - 1, Y^p - 1),
\]
we must show that
\[
\mathrm{Norm}(f) \equiv f(1, 1)^{p^2} \pmod{p^2 R}.
\]
For this, we view the norm as the composite of the norm maps in the $X$ and $Y$ variables separately:
\[
\begin{tikzcd}
R[X, Y]/(X^p - 1, Y^p - 1) \arrow[d, "N_X"'] \\
R[Y]/(Y^p - 1) \arrow[d, "N_Y"'] \\
R.
\end{tikzcd}
\]
Because the zero map $\Z/p\Z \to \mu_p$ is a generator over any $\F_p$-algebra, e.g.\ over $(R/pR)[Y]/(Y^p - 1)$, we have the congruence
\[
N_X(f(X, Y)) = f(1, Y)^p + p g(Y)
\]
with some $g(Y) \in R[Y]/(Y^p - 1)$. For any two elements $A, B$ in $R[Y]/(Y^p - 1)$ and an indeterminate $T$, the norm with respect to $Y$ of $A + TB$
\[
\begin{array}{ccc}
A + TB & \in & R[T][Y]/(Y^p - 1) \\
\Big\downarrow & & \Big\downarrow\rlap{$\scriptstyle N_Y$} \\
N_Y(A + TB) & \in & R[T],
\end{array}
\]
is a polynomial in $T$ of degree $p$, of the form
\[
N_Y(A + TB) = N_Y(A) + \sum_{i=1}^{p-1} T^i \Lambda_i(A, B) + T^p N_Y(B).
\]
But $\bmod p$ we have
\[
N_Y(A + TB) \equiv (A(1) + TB(1))^p \equiv A(1)^p + T^p B(1)^p \pmod p.
\]
Therefore equating coefficients yields
\[
\Lambda_i(A, B) \equiv 0 \pmod p \quad \text{for } i = 1, \dots, p-1,
\]
and therefore
\[
N_Y(A + TB) \equiv N_Y(A) \pmod{(pT, T^p)}.
\]
Applying this with $T = p$, we find
\[
N_Y(A + pB) \equiv N_Y(A) \pmod{p^2}.
\]
Therefore
\begin{align*}
\mathrm{Norm}(f(X, Y)) &= N_Y(N_X(f)) \\
&= N_Y(f(1, Y)^p + pg(Y)) \\
&\equiv N_Y((f(1, Y))^p) \pmod{p^2} \\
&\equiv (N_Y(f(1, Y)))^p \pmod{p^2} \\
&\equiv (f(1, 1)^p + pk)^p \pmod{p^2} \\
&\equiv f(1, 1)^{p^2} \pmod{p^2}.
\end{align*}

% ------------------------------------------------------------------------------
% Section 1.12
% ------------------------------------------------------------------------------
\section{Roots of unity}\label{sec:1.12}

\kmitem\label{km:1.12.1}
In this section, we will work out the general theory in the case of the multiplicative group $\mathbb{G}_m/\Z$, and $A = \Z/N\Z$. We denote by $\mu_N/\Z$ the group-scheme `\,$N$'th roots of unity', i.e., $\mu_N = \mathbb{G}_m[N]$.

\begin{lemma}\label{lem:1.12.2}
Over any scheme $S$, $(\mu_N)_S$ is the unique closed $S$-subgroup-scheme of $(\mathbb{G}_m)_S$ which is finite flat over $S$ of finite presentation and of rank $N$.
\end{lemma}

\begin{proof}
Visibly $(\mu_N)_S$ is such a closed $S$-subgroup-scheme of $(\mathbb{G}_m)_S$. If $G/S$ is any other, then $G$ is killed by $N$ (because it is of rank $N$). Therefore $G \subset \mathbb{G}_m[N] = \mu_N$, whence we have $G \le \mu_N$ as effective Cartier divisors in $(\mathbb{G}_m)_S/S$. Because both $G$ and $\mu_N$ have the same degree $N$, we may infer that $G = \mu_N$, as required.
\end{proof}

\begin{corollary}\label{cor:1.12.3}
If $A$ is any finite abelian ``abstract'' group of order $N$, we have a natural equality of functors
\[
A\text{\rm -Str}(\mathbb{G}_m/\Z) = A\text{\rm -Gen}(\mu_N/\Z),
\]
in particular
\[
(\Z/N\Z)\text{\rm -Str}(\mathbb{G}_m/\Z) = (\Z/N\Z)\text{\rm -Gen}(\mu_N/\Z).
\]
\end{corollary}

\begin{proof}
If $\phi$ is an $A$-structure on $\mathbb{G}_m$, then the $A$-subgroup it generates can be none other than $\mu_N$.
\end{proof}

\kmitem\label{km:1.12.4}
Henceforth we will consider only $A = \Z/N\Z$ (the natural choice since over any base scheme on which $N$ is invertible, $\mu_N$ is a twisted form of the constant group $\Z/N\Z$). We denote by $(\mu_N)^\times$ the closed subscheme of $\mu_N$ consisting of roots of unity of ``exact order $N$,'' i.e.,
\[
(\mu_N)^\times = (\Z/N\Z)\text{\rm -Str}(\mathbb{G}_m/\Z) = (\Z/N\Z)\text{\rm -Gen}(\mu_N/\Z).
\]
Thus for any ring $R$, we have:

\kmitem\label{km:1.12.5}
$(\mu_N)^\times(R) =$ the set of elements $\zeta \in R$ such that in the polynomial ring $R[T]$ we have the identity
\[
T^N - 1 = \prod_{a=1}^N (T - \zeta^a).
\]

\kmitem\label{km:1.12.6}
We will use interchangeably the expressions ``root of unity of `exact order $N$' in $R$'' and ``primitive $N$'th root of unity in $R$'' for elements of $(\mu_N)^\times(R)$.

\begin{theorem}\label{thm:1.12.7}
The scheme $(\mu_N)^\times$ of ``primitive $N$'th roots of unity'' is a regular one-dimensional scheme, which is finite and flat over $\Z$ of rank $\phi(N)$. Over $\Z[1/N]$, it is finite etale.
\end{theorem}

\begin{proof}
Suppose $N = N_1 N_2$ with $N_1$ and $N_2$ relatively prime. By \ref{km:1.10.5}, we have a canonical isomorphism
\[
(\mu_N)^\times \simeq (\mu_{N_1})^\times \underset{\Spec(\Z)}{\times} (\mu_{N_2})^\times.
\]
Therefore if the theorem is true for $N_1$ and $N_2$, then it is true for $N$.

This reduces us to treating the case when $N$ is a prime power $> 1$, say $N = p^n$. Over $\Z[1/p]$, we have a canonical isomorphism of $\Z[1/p]$-schemes
\[
(\mu_{p^n})^\times \xrightarrow{\ \sim\ } \underline{\Isom}(\Z/p^n\Z, \mu_{p^n}),
\]
which shows that, over $\Z[1/p]$, the scheme $(\mu_{p^n})^\times$ is indeed finite etale of rank $\phi(p^n)$.

The structural morphism
\[
\begin{tikzcd}
(\mu_{p^n})^\times \arrow[d] \\
\Spec(\Z)
\end{tikzcd}
\]
is surjective (because it is a finite morphism whose image, necessarily a closed subset of $\Spec(\Z)$, contains $\Spec(\Z[1/p])$). Because $\Spec(\Z)$ is one-dimensional, this surjectivity implies that there exists at least one closed point of $(\mu_{p^n})^\times$ lying over $p$ whose local ring has dimension $\ge 1$. In fact, there is at most one closed point of $(\mu_{p^n})^\times$ with characteristic $p$ (simply because for any field $k$ of characteristic $p$, $(\mu_{p^n})^\times(k) \subset \mu_{p^n}(k) = \{1\}$). Therefore $(\mu_{p^n})^\times$ has a unique closed point of characteristic $p$, namely the $\F_p$-valued point $\zeta = 1$. To show that $(\mu_{p^n})^\times$ is a regular one-dimensional scheme, it remains only to show that the local ring of $(\mu_{p^n})^\times$ at $\zeta = 1$ is regular and one-dimensional. As this local ring has dimension $\ge 1$, it suffices to show that its maximal ideal is generated by a single function. We claim that $\zeta - 1$ is such a function. Indeed, the global coordinate ring of $(\mu_{p^n})^\times$ is
\[
\Z[\zeta] \Big/ \Biggl( \parbox{0.78\linewidth}{\centering the relations obtained by equating coefficients in the identity $T^{p^n} - 1 = \prod_{a=1}^{p^n} (T - \zeta^a)$} \Biggr).
\]
We claim that if we divide this ring by the ideal $(\zeta - 1)$, the residue ring will be $\F_p$ itself. The residue ring is visibly
\[
\Z \Big/ \Biggl( \parbox{0.78\linewidth}{\centering the relations obtained by equating coefficients in $T^{p^n} - 1 = (T - 1)^{p^n}$} \Biggr).
\]
Comparing coefficients of $T$, we see that our ring is a quotient of $\Z/p^n\Z$. Comparing coefficients of $T^{p^{n-1}}$, and remembering that
\[
\ord_p \binom{p^n}{p^{n-1}} = 1,
\]
we see that our ring is a quotient of $\F_p$ itself. But our ring visibly maps onto $\F_p$, and hence it is $\F_p$. This concludes the proof that $(\mu_{p^n})^\times$ is a regular one-dimensional scheme.

To prove that $(\mu_{p^n})^\times$ is finite flat over $\Z$ of the asserted degree, it suffices to prove that it is finite and flat over $\Z$, for we know that over $\Z[1/p]$ it has the required degree. But $(\mu_{p^n})^\times$ is certainly finite over $\Spec(\Z)$, being a closed subscheme of $\mu_{p^n}$, and its flatness over $\Spec(\Z)$ results from the general fact that any finite morphism between regular schemes of the same dimension is automatically flat \cite[V, 3.6]{AK1}.
\end{proof}

\kmitem\label{km:1.12.8}
Let $\Phi_N(X) \in \Z[X]$ denote the $N$'th cyclotomic polynomial, i.e., the monic polynomial of degree $\phi(N)$ whose complex roots are the $\exp(2\pi i a/N)$, with $1 \le a \le N-1$ and $(a, N) = 1$. As is well-known, the ring
\[
\cO_N \xlongequal{dfn} \Z[X]/(\Phi_N(X))
\]
is precisely the ring of all algebraic integers in the field $\Q(\zeta_N)$ of $N$'th roots of unity.

\begin{theorem}\label{thm:1.12.9}
Over any ring $R$, an element $\zeta \in R$ is a root of the polynomial $\Phi_N$ if and only if $\zeta$ is a ``primitive $N$'th root of unity.'' Equivalently, we have an isomorphism of schemes
\[
\Spec(\cO_N) \xrightarrow{\ \sim\ } (\mu_N)^\times.
\]
\end{theorem}

\begin{proof}
We have a natural morphism
\[
\Spec(\cO_N) \longrightarrow (\mu_N)^\times,
\]
namely the $\cO_N$-valued point of $(\mu_N)^\times$ which is the element
\[
X \in \cO_N = \Z[X]/(\Phi_N(X)).
\]
(To see that this element $X$ actually lies in $(\mu_N)^\times(\cO_N)$, i.e., that it satisfies the identity
\[
T^N - 1 = \prod_{a=1}^N (T - X^a),
\]
it suffices to do so after any injective extension of scalars $\cO_N \hookrightarrow K$. We use the embedding $\cO_N \hookrightarrow \C$ given by $X \mapsto \exp(2\pi i/N)$, which makes this identity obvious.)

This morphism is certainly an isomorphism over $\Z[1/N]$, (because both source and target are finite etale of the same rank $\phi(N)$ over $\Z[1/N]$ it suffices to check that for any algebraically closed field $k$ in which $N$ is invertible, the map induces an isomorphism on $k$-valued points. By \ref{lem:1.4.4}~(3) this amounts to the statement that the roots of $\Phi_N$ in such a field $k$ are precisely those $N$'th roots of unity in $k$ which are not $M$'th roots of unity for any $1 \le M < N$, which is standard).

We next define a morphism
\[
(\mu_N)^\times \longrightarrow \Spec(\cO_N),
\]
which carries the universal primitive $N$'th root of unity $\zeta$ to the universal solution $X$ of $\Phi_N(X) = 0$. To do this, we must show that the element $\zeta$ in the coordinate ring $B$ of $(\mu_N)^\times$
\[
\small \zeta \in B = \Z[\zeta] \Big/ \Biggl( \parbox{0.70\linewidth}{\centering the relations obtained by equating coefficients in\\[1pt] $T^N - 1 = \prod_{a=1}^N (T - \zeta^a)$} \Biggr)
\]
is a root of $\Phi_N$ in that same ring, i.e., that
\[
\Phi_N(\zeta) = 0 \quad \text{in } B.
\]
But $B$ is flat over $\Z$, so in particular $B \subset B \underset{\Z}{\otimes} \Z[1/N]$. Therefore it suffices to verify that
\[
\Phi_N(\zeta) = 0 \quad \text{in } B[1/N].
\]
But this vanishing amounts precisely to the just-proven statement that the natural morphism
\[
\Spec(\cO_N) \longrightarrow (\mu_N)^\times
\]
is an isomorphism over $\Z[1/N]$.

To verify that each of the two composite maps
\begin{gather*}
\Spec(\cO_N) \longrightarrow (\mu_N)^\times \longrightarrow \Spec(\cO_N) \\
(\mu_N)^\times \longrightarrow \Spec(\cO_N) \longrightarrow (\mu_N)^\times
\end{gather*}
is the identity, it again suffices to check over $\Z[1/N]$ (since both $\cO_N$ and $(\mu_N)^\times$ are flat over $\Z$), where this is obvious.
\end{proof}

\begin{remark}\label{rem:1.12.10}
The arguments given in this section are in many ways typical of the arguments which will be used later in discussing moduli of elliptic curves. This perhaps justifies giving them in such detail.
\end{remark}

% ------------------------------------------------------------------------------
% Section 1.13
% ------------------------------------------------------------------------------
\section{Some open problems}\label{sec:1.13}

In this section we would like to mention explicitly some open problems related to the ideas developed in this chapter.

\begin{enumerate}[\bf 1.]
\item Let $S$ be an arbitrary scheme, $Z/S$ a finite locally free $S$-scheme of rank $N \ge 2$. Is it true that locally f.p.p.f.\ on $S$, $Z/S$ always admits a full set of sections? (It is of course true if $Z/S$ is embeddable in a smooth curve $C/S$.)

\item Let $S$ be an arbitrary scheme, $Z_1/S$ and $Z_2/S$ two finite locally free $S$-schemes of the same rank $N \ge 1$. Let
\[
\phi \colon Z_1 \longrightarrow Z_2
\]
be an $S$-morphism. We say that $\phi$ is a \emph{$\times$-morphism} if for every affine $S$-scheme $\Spec(R)$, and for every function $f$ on $(Z_2)_R = \Spec(B_2)$ we have an equality of characteristic polynomials in $R[T]$
\[
\det(T - f) = \det(T - \phi^*(f)).
\]
Clearly the composition of two $\times$-morphisms is another. When $Z_1$ is the disjoint union of $N$ copies of $S$, such a $\times$-morphism is exactly what we have called a full set of sections of $Z_2/S$.

Now suppose that $G_1$ and $G_2$ are finite locally free commutative group-schemes over $S$, of the same rank $N$. We say that a homomorphism of $S$-group-schemes
\[
\phi \colon G_1 \longrightarrow G_2
\]
is a \emph{$\times$-homomorphism} if it is a $\times$-morphism of the underlying $S$-schemes. Clearly the composition of two $\times$-homomorphisms is another one. So for given $N \ge 1$, we can form a category by taking as objects the finite locally free commutative $S$-group-schemes of rank $N$, and as morphisms the set
\[
\Hom_{S\text{\rm -gp}}^\times(G_1, G_2)
\]
of all $\times$-homomorphisms. If $G_1$ is a constant abelian group $A$ of order $N$, then we have
\[
\Hom_S^\times((A)_S, G_2) = \text{the set of $A$-generators of $G_2/S$}.
\]
Given $G_1, G_2$, let us say that $G_2$ is a ``degeneration'' of $G_1$ if, after some f.p.p.f.\ base-change, $T \to S$, the set
\[
\Hom_T^\times(G_1, G_2),
\]
which we might call the ``$G_1$-generators of $G_2$,'' is non-empty. Is there an interesting theory of ``degeneration''? Are there cases ``in nature'' where one wants to consider moduli problems based on the notion of a ``$G_1$-generator'' when $G_1$ is not constant?
\end{enumerate}
